% Self-contained source generated from the July 18, 2026 Pass 7 project.
% All project-local TeX inputs are inlined below.
\documentclass[11pt,reqno]{scrartcl}

\usepackage[margin=1in]{geometry}
\usepackage[T1]{fontenc}
\usepackage{lmodern}
\usepackage{microtype}
\usepackage{setspace}
\usepackage{amsmath,amssymb,amsthm,mathtools,mathrsfs}
\usepackage{booktabs,tabularx,array}
\usepackage{enumitem}
\usepackage{xcolor}
\usepackage[colorlinks=true,linkcolor=blue!60!black,citecolor=blue!60!black,urlcolor=blue!60!black]{hyperref}
\hypersetup{
  pdftitle={Formalizing Buddhist (Meta)physics: Toward a Universal Model},
  pdfauthor={Alexander Yiannopoulos},
  pdfsubject={Bipartite Krein transition algebra, compact symmetry compression, reference-state geometry, projector topology, effective dynamics, and passive response},
  pdfkeywords={Krein space, closed operator, graded algebra, Haar projector, reference fidelity, projection bundle, Chern-Weil theory, gauge curvature, Schur complement, Herglotz function, Schur function, response realization}
}
\usepackage{aliascnt}
\usepackage[nameinlink,capitalise,noabbrev]{cleveref}
\usepackage{import}
\usepackage{scrlayer-scrpage}


% ===== BEGIN INLINED FILE: macros.tex =====
% Core notation for Toward a Universal Model
\newcommand{\CC}{\mathbb C}
\newcommand{\RR}{\mathbb R}
\newcommand{\ZZ}{\mathbb Z}
\newcommand{\NN}{\mathbb N}

\newcommand{\scrK}{\mathscr K}
\newcommand{\scrHp}{\mathscr H_{+}}
\newcommand{\scrHm}{\mathscr H_{-}}
\newcommand{\fundJ}{\mathfrak J}
\newcommand{\Pp}{\Pi_{+}}
\newcommand{\Pm}{\Pi_{-}}
\newcommand{\idop}{\mathbf 1}

\newcommand{\Aop}{\mathsf A}
\newcommand{\Vop}{\mathsf V}
\newcommand{\DB}{\mathbb D}
\newcommand{\KB}{\mathbb K}
\newcommand{\DeltaB}{\Delta}
\newcommand{\Cp}{\mathscr C_{+}}
\newcommand{\Cm}{\mathscr C_{-}}

\newcommand{\ket}[1]{\lvert #1\rangle}
\newcommand{\bra}[1]{\langle #1\rvert}
\newcommand{\braket}[2]{\langle #1\mid #2\rangle}
\newcommand{\inner}[2]{\langle #1,#2\rangle}
\newcommand{\kInner}[2]{[#1,#2]_{\fundJ}}

\DeclareMathOperator{\Dom}{Dom}
\DeclareMathOperator{\Ran}{Ran}
\DeclareMathOperator{\Ker}{Ker}
\DeclareMathOperator{\spec}{spec}
\DeclareMathOperator{\diag}{diag}
\DeclareMathOperator{\ind}{ind}
\DeclareMathOperator{\Graph}{Graph}
\DeclareMathOperator{\Tr}{Tr}
\DeclareMathOperator{\coker}{coker}
\DeclareMathOperator{\Span}{span}

\newcommand{\restr}[2]{#1\!\upharpoonright_{#2}}
\newcommand{\dsum}{\mathbin{[\dotplus]}}

% Polar geometry and cohomology
\newcommand{\polu}{u_B}
\newcommand{\suppP}{\pi_{+,B}}
\newcommand{\suppM}{\pi_{-,B}}
\newcommand{\SuppB}{\Pi_B}
\newcommand{\RB}{\mathcal R_B}
\newcommand{\IB}{\mathcal I_B}
\newcommand{\Harm}{\mathscr H}

% Reserved notation for later realizations
\newcommand{\scalarpot}[1]{U\!\left(#1\right)}
\newcommand{\gaugeA}{\mathcal A}
\newcommand{\gaugeF}{\mathcal F}

% Boundary response layer
\newcommand{\scrX}{\mathscr X}
\newcommand{\scrE}{\mathscr E}
\newcommand{\Mresp}{M_{\Gamma}}
\newcommand{\Sresp}{\mathscr S_{\Gamma}}
\newcommand{\KM}{\mathcal K_{M}}
\newcommand{\KS}{\mathcal K_{\mathscr S}}
\newcommand{\Cplus}{\mathbb C_{+}}
\newcommand{\Hmin}{\mathscr X_{\min}}

% Reference-state geometry
\newcommand{\vref}[1]{v_{#1}}
\newcommand{\Zref}[1]{Z_{#1}}
\DeclareMathOperator{\rank}{rank}

% Integration pass
% ===== END INLINED FILE: macros.tex =====


\setstretch{1.15}
\setlength{\emergencystretch}{3em}
\setkomafont{section}{\Large\bfseries}
\setkomafont{subsection}{\large\bfseries}
\setkomafont{disposition}{\normalfont}
\DeclareTOCStyleEntry[numwidth=3.5em]{tocline}{subsection}

\clearpairofpagestyles
\ihead{\textit{Toward a Universal Model}}
\ohead{}
\cfoot{\pagemark}

\theoremstyle{plain}
\newtheorem{theorem}{Theorem}[section]
\newaliascnt{proposition}{theorem}
\newtheorem{proposition}[proposition]{Proposition}
\aliascntresetthe{proposition}
\newaliascnt{lemma}{theorem}
\newtheorem{lemma}[lemma]{Lemma}
\aliascntresetthe{lemma}
\newaliascnt{corollary}{theorem}
\newtheorem{corollary}[corollary]{Corollary}
\aliascntresetthe{corollary}
\theoremstyle{definition}
\newaliascnt{definition}{theorem}
\newtheorem{definition}[definition]{Definition}
\aliascntresetthe{definition}
\newaliascnt{assumption}{theorem}
\newtheorem{assumption}[assumption]{Assumption}
\aliascntresetthe{assumption}
\newaliascnt{example}{theorem}
\newtheorem{example}[example]{Example}
\aliascntresetthe{example}
\theoremstyle{remark}
\newaliascnt{remark}{theorem}
\newtheorem{remark}[remark]{Remark}
\aliascntresetthe{remark}

\crefname{theorem}{theorem}{theorems}
\crefname{proposition}{proposition}{propositions}
\crefname{lemma}{lemma}{lemmas}
\crefname{corollary}{corollary}{corollaries}
\crefname{definition}{definition}{definitions}
\crefname{assumption}{assumption}{assumptions}
\crefname{example}{example}{examples}
\crefname{section}{section}{sections}
\crefname{appendix}{appendix}{appendices}

\title{\textbf{Formalizing Buddhist (Meta)physics: Toward a Universal Model}}
\author{Alexander Yiannopoulos\\[2pt]
\small A$\lvert\Omega\rangle$ Institute, New Orleans, Louisiana, USA\\
\small\texttt{ayianno@ao.institute}}
\date{July 18, 2026; revised November 12, 2026}

\begin{document}
\maketitle

\begin{abstract}
\noindent
Let \(\scrK=\scrHp\dsum\scrHm\) be a bipartite Krein space and let
\(B:\Dom B\subseteq\scrHp\to\scrHm\) be densely defined and closed.  We
construct its nilpotent oriented transitions, Hilbert/Hodge and
Krein/rotational odd branches, positive partner Laplacian, polar geometry,
two-term cohomology, and paired nonzero spectra.  When \(B\) is Fredholm,
its zero-mode imbalance gives \(\ind B\); under the stated trace-class
heat-kernel hypotheses, the same index is the graded heat supertrace.  A
norm-\(C^2\) field of orthogonal projections produces bounded tangent channels
\(B_\mu\) and the Grassmann curvature
\(\gaugeF_{\mu\nu}^{\Pi}=B_\mu^*B_\nu-B_\nu^*B_\mu\).  Under the additional
smooth finite-rank hypothesis, its trace polynomials give the Chern--Weil
forms.  A bounded channel coupling
between self-adjoint diagonal generators yields the sign-sensitive
Schur--Feshbach operator
\(L_{\mathrm{eff}}^{(\varepsilon)}(z)=L_+-\varepsilon
B^*(L_--z)^{-1}B\).  A bounded port into a self-adjoint realization gives
an operator-valued Herglotz response, a Schur Cayley transform, and positive
Nevanlinna and de Branges--Rovnyak kernels; minimal realizations are unique up
to unitary equivalence.  The response-minimal carrier separates
dynamically reached final defects from permanently unobserved ambient
directions.  A symmetry-resolved response theorem describes departure
from a port, and an explicit three-state example preserves a ground
while exhibiting complete departure from a two-dimensional preparation
range.  Gaussian functional calculus regularizes
singular boundary inputs without changing the internal spectrum.  A supplied
ground state also gives the exact reference-fidelity decomposition
\(E_{\mathrm{exc}}(\psi)=Z_\psi(1-|\braket0\psi|^2)\) off the reference ray,
with zero excitation on that ray.  Scalar,
translation, oscillator, and quasiperiodic channels exhibit finite,
continuous, positive-integer, and compact discrete spectral types.  For the
oscillator channel \(B=a^*\), the Mellin transform of the heat trace of the
squared positive return,
\(\operatorname{Tr}e^{-\pi t(B^*B)^2}\), equals
\(\pi^{-s/2}\Gamma(s/2)\zeta(s)\) for \(\operatorname{Re}s>1\).  An
explicit scope ledger separates this established local architecture from
extensions requiring additional realization data.
\end{abstract}

\medskip
\noindent\textbf{Keywords:} Krein space; closed operator; graded algebra;
two-term Hilbert complex; compact-group averaging; reference fidelity;
projection bundle; Chern--Weil theory; gauge curvature; Schur complement;
Herglotz response; Schur transfer; minimal realization; heat-kernel smoothing.

\clearpage
\setcounter{tocdepth}{2}
\tableofcontents

\clearpage

% ===== BEGIN INLINED FILE: sections/00_introduction.tex =====
\section{Introduction and scope}
\label{sec:introduction}

This paper develops a common operator architecture for systems with two
complementary sectors, an indefinite grading, and a positive compressed
response.  The starting datum is deliberately small: a bipartite Krein space
and one densely defined closed channel between its two Hilbert sectors.  From
that datum one obtains two nilpotent directional maps, two canonical odd
branches, a positive partner Laplacian, a supportwise polar geometry, a
two-term cohomology, and, under Fredholm hypotheses, an index.  When the
selected sector varies over a
parameter space, the same algebra yields the Grassmann connection and its
curvature.  When the complementary sector is eliminated through a resolvent,
the same closed excursion yields a Schur--Feshbach self-energy.  Finally, a
bounded port into a self-adjoint realization produces a Herglotz response, a
Schur transfer function, and positive reproducing kernels.

The adjective \emph{universal} refers to this architecture, not to a fixed
spectrum or a fixed physical realization.  The channel operator determines
the spectral type.  The scalar, translation, oscillator, and quasiperiodic
examples in \cref{sec:examples} therefore instantiate the same algebra with,
respectively, finite, continuous, positive-integer, and compact discrete
spectra.

Most component tools used below--polar decomposition, two-term Hilbert
complexes, Grassmann connections, Schur complements, and Nevanlinna kernels--
are classical.  The contribution of the present paper is their classification
and synthesis around one oriented channel pair, including the exact sign
dichotomy between the Hilbert and Krein branches and the recurrence of the
same closed-return pattern in curvature, effective dynamics, and passive
response.

\subsection{Primitive data and derived structures}

The fixed-channel layer begins with a fundamental decomposition
\(\scrK=\scrHp\dsum\scrHm\), its symmetry \(\fundJ\), and one densely
defined closed channel
\(B:\Dom B\subseteq\scrHp\to\scrHm\).  The associated differential and
codifferential are denoted by \(d_B\) and \(d_B^*\), and
\begin{equation}
  \Aop_B:=-d_B,
  \qquad
  \Vop_B:=d_B^*.
  \label{eq:intro-AV}
\end{equation}
The letters \(\Aop\) and \(\Vop\) label the two orientations of the
transition channel.  Their Hilbert/Hodge and Krein/rotational branches are
\begin{equation}
  \DB_B=\Vop_B-\Aop_B,
  \qquad
  \KB_B=\Vop_B+\Aop_B=\fundJ\DB_B,
\end{equation}
with
\begin{equation}
  \DB_B^2=\DeltaB_B,
  \qquad
  \KB_B^2=-\DeltaB_B,
  \qquad
  \DeltaB_B=\diag(B^*B,BB^*)\ge0.
\end{equation}
All unbounded products use their natural composition domains, collected in
\hyperref[app:domains]{Appendix~\ref*{app:domains}}.  Independently, a compact unitary symmetry may be
averaged to the invariant projection \(\Pi_G\); this projection is kept
distinct from the fixed sector projections, the support of \(B\), and a
moving projection field.

\subsection{One closed excursion in three forms}

The same algebra appears in three mathematically distinct settings.

First, for a norm-\(C^2\) orthogonal projection field \(\Pi(x)\), the tangent
channel
\begin{equation}
  B_\mu=\Pi^\perp(\partial_\mu\Pi)\Pi
\end{equation}
produces the projector curvature
\begin{equation}
  \gaugeF_{\mu\nu}^{\Pi}
  =B_\mu^*B_\nu-B_\nu^*B_\mu.
  \label{eq:intro-curvature}
\end{equation}
Thus curvature is the antisymmetric difference of two mixed returns.

Second, let \(L_-\) be self-adjoint, let
\(B\in\mathcal B(\scrHp,\scrHm)\), and take \(z\in\rho(L_-)\).  The
resolvent return
\begin{equation}
  \Sigma_B(z)=B^*(L_--z)^{-1}B
  \label{eq:intro-self-energy}
\end{equation}
produces the sign-sensitive effective operator
\begin{equation}
  L_{\mathrm{eff}}^{(\varepsilon)}(z)
  =L_+-\varepsilon\Sigma_B(z),
  \qquad \varepsilon\in\{+1,-1\}.
  \label{eq:intro-effective}
\end{equation}

Third, for a self-adjoint response generator \(H\) and bounded port
\(\Gamma\), the compressed resolvent
\begin{equation}
  M_\Gamma(z)=\Gamma^*(H-z)^{-1}\Gamma
  \label{eq:intro-response}
\end{equation}
is Herglotz--Nevanlinna on the upper half-plane.  Its Cayley transform is
Schur and its defect kernel is positive.  These three constructions are not
identified with one another; rather, they are generated by the same
orientation-and-return pattern under different hypotheses.

\subsection{Reference states and topology}

A normalized vector \(\ket0\in\Dom B\subseteq\scrHp\) may be chosen as a
distinguished
reference.  Its one-step image is
\begin{equation}
  \ket\Omega:=\Aop_B\ket0,
  \qquad
  \Aop_B\ket\Omega=0.
  \label{eq:intro-reference-image}
\end{equation}
No ground-state or cyclicity property is included in this definition.  When a
self-adjoint realization additionally makes \(\ket0\) a ground state, the
reference-state theorem of \cref{sec:reference-state} gives an exact
factorization of excitation energy by the fidelity with the reference ray.

For a strongly continuous unitary action of a compact group, the Haar
average gives the exact orthogonal projection onto invariant vectors; see
\cref{sec:symmetry-averaging}.  This construction supplies invariant
compression but does not identify that compression with the selected sector
unless a realization proves the identification.

For a finite-rank smooth projection bundle, the curvature in
\eqref{eq:intro-curvature} also generates the standard Chern--Weil forms.  In
rank one, the resulting flux through a closed oriented two-cycle is integral.
This supplies the topological consequence of the projector construction
without imposing a quantization condition on the holonomy of every closed
loop.

\subsection{Logical organization}

The paper is divided into four layers.
\begin{enumerate}[label=(\roman*)]
  \item \emph{Fixed-channel algebra:} closed operators, polar geometry,
  cohomology, spectral pairing, and index.
  \item \emph{Variable compression:} projector curvature and
  characteristic classes.
  \item \emph{Effective and boundary dynamics:} Schur--Feshbach reduction,
  Herglotz response, Schur kernels, and minimal realization.
  \item \emph{Realizations and open extensions:} canonical examples and a
  precise ledger of which further constructions require additional data.
\end{enumerate}
The exact hypotheses of each result are local to its theorem.  No example is
used to prove the abstract algebra, and no spectral identity from one
realization is promoted to a universal spectrum.
% ===== END INLINED FILE: sections/00_introduction.tex =====


% ===== BEGIN INLINED FILE: sections/01_setup.tex =====
\section{Bipartite Krein geometry}
\label{sec:setup}

\subsection{Fundamental decomposition}

Let \(\scrHp\) and \(\scrHm\) be complex separable Hilbert spaces. Their inner
products are written with the same brackets and are taken to be linear in the
second argument. Set
\begin{equation}
  \scrK:=\scrHp\oplus\scrHm
\end{equation}
with the direct-sum Hilbert product
\begin{equation}
  \inner{(x_+,x_-)}{(y_+,y_-)}
  :=\inner{x_+}{y_+}_{+}+\inner{x_-}{y_-}_{-}.
\end{equation}
Define the bounded self-adjoint involution
\begin{equation}
  \fundJ:=
  \begin{pmatrix}
    \idop_{+}&0\\
    0&-\idop_{-}
  \end{pmatrix},
  \qquad
  \fundJ^*=\fundJ,
  \qquad
  \fundJ^2=\idop_{\scrK}.
  \label{eq:J}
\end{equation}
The associated indefinite form is
\begin{equation}
  \kInner{\xi}{\eta}:=\inner{\xi}{\fundJ\eta},
  \qquad \xi,\eta\in\scrK.
  \label{eq:Krein-form}
\end{equation}
Thus, after the canonical sector embeddings into \(\scrK\),
\begin{equation}
  \kInner{(x_+,0)}{(x_+,0)}=\|x_+\|^2,
  \qquad
  \kInner{(0,x_-)}{(0,x_-)}=-\|x_-\|^2.
\end{equation}
and the corresponding fundamental decomposition is written
\begin{equation}
  \boxed{\scrK=\scrHp\dsum\scrHm.}
  \label{eq:fundamental-decomposition}
\end{equation}
This is a Krein space in the standard sense \cite{Bognar,AzizovIokhvidov}.
The Hilbert topology is used for closure, domains, norms, and ordinary
adjoints; the indefinite form records the sector grading and defines the
Krein adjoint.

The fixed sector projections are
\begin{equation}
  \Pp:=
  \begin{pmatrix}\idop_+&0\\0&0\end{pmatrix},
  \qquad
  \Pm:=
  \begin{pmatrix}0&0\\0&\idop_-\end{pmatrix},
  \label{eq:sector-projections}
\end{equation}
so that
\begin{equation}
  \Pp+\Pm=\idop_{\scrK},
  \qquad
  \Pp\Pm=0,
  \qquad
  \fundJ=\Pp-\Pm.
\end{equation}

\subsection{Adjoints and grading}

For a densely defined operator \(X\) on \(\scrK\), its Hilbert adjoint is
written \(X^*\). Whenever the expression is defined as a closed operator, its
Krein adjoint is
\begin{equation}
  X^\sharp:=\fundJ X^*\fundJ.
  \label{eq:Krein-adjoint}
\end{equation}
The domain is
\begin{equation}
  \Dom(X^\sharp)=\fundJ\Dom(X^*).
\end{equation}

\begin{definition}[Grading-odd operator]
A densely defined operator \(X\) on \(\scrK\) is \emph{grading-odd} when
\(\fundJ\Dom X=\Dom X\) and
\begin{equation}
  \fundJ X\xi=-X\fundJ\xi,
  \qquad \xi\in\Dom X.
  \label{eq:odd}
\end{equation}
Equivalently, its diagonal sector blocks vanish on their natural domains:
\begin{equation}
  \Pp X\Pp=0,
  \qquad
  \Pm X\Pm=0.
\end{equation}
\end{definition}

The \(\fundJ\)-invariance of the domain is essential for unbounded operators:
it permits the domain itself to split into its two sector components.

\begin{lemma}[Odd-adjoint conversion]
\label{lem:odd-adjoint}
Let \(X\) be densely defined, closed, and grading-odd. Then
\begin{equation}
  X^\sharp=X
  \quad\Longleftrightarrow\quad
  X^*=-X.
  \label{eq:odd-adjoint-conversion}
\end{equation}
Likewise,
\begin{equation}
  X^\sharp=-X
  \quad\Longleftrightarrow\quad
  X^*=X.
  \label{eq:odd-adjoint-conversion-2}
\end{equation}
\end{lemma}

\begin{proof}
If \(X^\sharp=X\), then \(X^*=\fundJ X\fundJ\), including equality of
domains. Oddness gives \(\fundJ X\fundJ=-X\), hence \(X^*=-X\). The reverse
implication is the same calculation in the opposite order. Replacing
\(X^\sharp=X\) by \(X^\sharp=-X\) gives the second equivalence.
\end{proof}

\begin{remark}
For an odd operator, Krein self-adjointness is therefore equivalent to
Hilbert skew-adjointness. Krein skew-adjointness is equivalent to Hilbert
self-adjointness. The sign is structural and will distinguish the two
canonical odd branches associated with one channel.
\end{remark}

\subsection{Equal-carrier realization}

The universal construction does not require \(\scrHp\) and \(\scrHm\) to be
isomorphic. If they are specified as copies of a common Hilbert space
\(\mathscr H\), however, then
\begin{equation}
  \scrK\simeq\CC^{1,1}\otimes\mathscr H
  \simeq\mathscr H\oplus\mathscr H,
\end{equation}
with
\begin{equation}
  \fundJ=\sigma_z\otimes\idop_{\mathscr H}.
\end{equation}
This equal-carrier representation is a specialization of the direct-sum
geometry, not a replacement for it.
% ===== END INLINED FILE: sections/01_setup.tex =====


% ===== BEGIN INLINED FILE: sections/02_classification.tex =====
\section{Closed channels and canonical odd branches}
\label{sec:classification}

\subsection{The two-dimensional sector factor}

Before introducing an operator-valued channel, consider
\begin{equation}
  M=
  \begin{pmatrix}
    0&a\\
    b&0
  \end{pmatrix},
  \qquad a,b\in\CC,
\end{equation}
on \(\CC^{1,1}\) with fundamental symmetry \(\sigma_z\).

\begin{proposition}[Canonical sector matrix]
\label{prop:canonical-sector}
The matrix \(M\) is Krein-self-adjoint if and only if
\begin{equation}
  b=-\overline a.
\end{equation}
If \(a\neq0\), a diagonal \(\sigma_z\)-unitary phase transformation carries
\(M\) to
\begin{equation}
  |a|\,i\sigma_y
  =|a|
  \begin{pmatrix}
    0&1\\
    -1&0
  \end{pmatrix}.
  \label{eq:i-sigma-y}
\end{equation}
\end{proposition}

\begin{proof}
A direct calculation gives
\begin{equation}
  M^\sharp
  =\sigma_zM^*\sigma_z
  =
  \begin{pmatrix}
    0&-\overline b\\
    -\overline a&0
  \end{pmatrix}.
\end{equation}
Thus \(M^\sharp=M\) is equivalent to \(b=-\overline a\). Write
\(a=|a|e^{i\theta}\) and define
\begin{equation}
  U_\theta:=\diag(e^{-i\theta/2},e^{i\theta/2}).
\end{equation}
Then \(U_\theta^*\sigma_zU_\theta=\sigma_z\) and
\(U_\theta M U_\theta^{-1}=|a|i\sigma_y\).
\end{proof}

\subsection{A closed channel and its two-term complex}

\begin{assumption}[Closed transition channel]
\label{ass:B}
Let
\begin{equation}
  B:\Dom B\subseteq\scrHp\longrightarrow\scrHm
  \label{eq:B}
\end{equation}
be densely defined and closed. We use the standard graph, closure, and adjoint
conventions for unbounded operators \cite{Kato,ReedSimon}. Its Hilbert adjoint is
\begin{equation}
  B^*:\Dom B^*\subseteq\scrHm\longrightarrow\scrHp.
\end{equation}
\end{assumption}

Define the two-term differential
\begin{equation}
  d_B:=
  \begin{pmatrix}
    0&0\\
    B&0
  \end{pmatrix},
  \qquad
  \Dom d_B:=\Dom B\oplus\scrHm.
  \label{eq:dB}
\end{equation}
Its Hilbert adjoint is
\begin{equation}
  d_B^*=
  \begin{pmatrix}
    0&B^*\\
    0&0
  \end{pmatrix},
  \qquad
  \Dom d_B^*:=\scrHp\oplus\Dom B^*.
  \label{eq:dBstar}
\end{equation}

\begin{definition}[Oriented transition pair]
\label{def:AV}
The two oriented transitions associated with \(B\) are
\begin{equation}
  \boxed{
  \Aop_B:=-d_B
  =
  \begin{pmatrix}
    0&0\\
    -B&0
  \end{pmatrix},
  \qquad
  \Vop_B:=d_B^*
  =
  \begin{pmatrix}
    0&B^*\\
    0&0
  \end{pmatrix}.}
  \label{eq:AV}
\end{equation}
The letters \(\Aop\) and \(\Vop\) label the two orientations of the
transition channel.
\end{definition}

\begin{proposition}[Closedness, adjoints, and nilpotency]
\label{prop:differential-basic}
The operators \(d_B\), \(d_B^*\), \(\Aop_B\), and \(\Vop_B\) are densely
defined and closed. Moreover,
\begin{equation}
  d_B^2=0,
  \qquad
  (d_B^*)^2=0,
  \label{eq:d-nilpotent}
\end{equation}
with
\begin{equation}
  \Dom(d_B^2)=\Dom d_B,
  \qquad
  \Dom((d_B^*)^2)=\Dom d_B^*.
\end{equation}
The Hilbert and Krein adjoints of the oriented pair are
\begin{align}
  \Aop_B^*&=-\Vop_B,
  &\Vop_B^*&=-\Aop_B,
  \label{eq:AV-Hilbert-adjoint}\\
  \Aop_B^\sharp&=\Vop_B,
  &\Vop_B^\sharp&=\Aop_B.
  \label{eq:AV-Krein-adjoint}
\end{align}
\end{proposition}

\begin{proof}
Closedness of \(d_B\) is equivalent to closedness of \(B\), because its graph
is the direct product of \(\Graph(B)\) with the unrestricted second-sector
coordinate. The adjoint calculation gives \eqref{eq:dBstar}, hence \(d_B^*\)
is closed. The remaining closedness claims follow from multiplication by
\(-1\).

For \((x,y)\in\Dom d_B\), one has \(d_B(x,y)=(0,Bx)\). This output lies again
in \(\Dom d_B\), because its first component is zero; therefore
\(d_B^2(x,y)=0\). The codifferential case is identical. Finally,
\(\Aop_B^*=-d_B^*=-\Vop_B\). Conjugation by \(\fundJ\) changes the sign of an
off-diagonal block, giving \(\Aop_B^\sharp=\Vop_B\); the other identities
follow symmetrically.
\end{proof}

\subsection{Classification of odd Krein-self-adjoint generators}

The oriented pair combines into the Krein branch
\begin{equation}
  \KB_B:=\Aop_B+\Vop_B
  =
  \begin{pmatrix}
    0&B^*\\
    -B&0
  \end{pmatrix},
  \qquad
  \Dom\KB_B=\Dom B\oplus\Dom B^*.
  \label{eq:KB}
\end{equation}

\begin{theorem}[Odd Krein-self-adjoint classification]
\label{thm:classification}
Let \(X\) be a densely defined, closed, grading-odd operator on
\(\scrK\) and suppose \(X^\sharp=X\). Then
there exists a unique densely defined closed operator
\(B:\Dom B\subseteq\scrHp\to\scrHm\) such that
\begin{equation}
  X=\KB_B
  =
  \begin{pmatrix}
    0&B^*\\
    -B&0
  \end{pmatrix}.
  \label{eq:classification-form}
\end{equation}
Conversely, every operator \(\KB_B\) defined by a densely defined closed
channel \(B\) is closed, grading-odd, Hilbert-skew-adjoint, and
Krein-self-adjoint.
\end{theorem}

\begin{proof}
Because \(\Dom X\) is invariant under \(\fundJ\), it is invariant under
\(\Pp=\tfrac12(\idop+\fundJ)\) and
\(\Pm=\tfrac12(\idop-\fundJ)\). Hence its domain decomposes into dense sector
domains. Oddness gives a block form
\begin{equation}
  X=
  \begin{pmatrix}
    0&C\\
    D&0
  \end{pmatrix},
\end{equation}
where
\(D:\Dom D\subseteq\scrHp\to\scrHm\) and
\(C:\Dom C\subseteq\scrHm\to\scrHp\). Closedness of \(X\) implies
closedness of both restrictions. By \cref{lem:odd-adjoint},
\(X^*=-X\). The block adjoint is
\begin{equation}
  X^*=
  \begin{pmatrix}
    0&D^*\\
    C^*&0
  \end{pmatrix},
\end{equation}
so \(D^*=-C\) and \(C^*=-D\). Set \(B:=-D\). Then \(C=B^*\), giving
\eqref{eq:classification-form}. Uniqueness is immediate from the lower-left
block.

Conversely, the direct block calculation gives \(\KB_B^*=-\KB_B\).
Closedness follows from closedness of \(B\) and \(B^*\), and oddness is
manifest. The odd-adjoint conversion then yields \(\KB_B^\sharp=\KB_B\).
\end{proof}

\subsection{Hilbert and Krein branches}

The same channel \(B\) determines a second canonical odd branch,
\begin{equation}
  \DB_B:=\Vop_B-\Aop_B
  =d_B+d_B^*
  =
  \begin{pmatrix}
    0&B^*\\
    B&0
  \end{pmatrix},
  \qquad
  \Dom\DB_B=\Dom B\oplus\Dom B^*.
  \label{eq:DB}
\end{equation}
Define the partner Laplacian
\begin{equation}
  \DeltaB_B:=
  \begin{pmatrix}
    B^*B&0\\
    0&BB^*
  \end{pmatrix},
  \label{eq:DeltaB}
\end{equation}
with
\begin{equation}
  \Dom\DeltaB_B
  :=\Dom(B^*B)\oplus\Dom(BB^*).
\end{equation}
The operators \(B^*B\) and \(BB^*\) are positive and self-adjoint \cite{Kato,ReedSimon}; hence so
is \(\DeltaB_B\).

\begin{theorem}[Adjoint-sign dichotomy]
\label{thm:dichotomy}
The Hilbert/Hodge branch \(\DB_B\) and the Krein/rotational branch
\(\KB_B\) satisfy
\begin{align}
  \DB_B^*&=\DB_B,
  &\DB_B^\sharp&=-\DB_B,
  &\DB_B^2&=+\DeltaB_B,
  \label{eq:D-branch}\\
  \KB_B^*&=-\KB_B,
  &\KB_B^\sharp&=\KB_B,
  &\KB_B^2&=-\DeltaB_B.
  \label{eq:K-branch}
\end{align}
The two branches are related exactly by
\begin{equation}
  \boxed{\KB_B=\fundJ\DB_B.}
  \label{eq:KJD}
\end{equation}
\end{theorem}

\begin{proof}
The adjoint statements follow from the block matrices and
\cref{lem:odd-adjoint}. Direct multiplication gives the square identities on
\(\Dom\DeltaB_B\). Finally,
\begin{equation}
  \fundJ\DB_B
  =
  \begin{pmatrix}
    \idop_+&0\\0&-\idop_-
  \end{pmatrix}
  \begin{pmatrix}
    0&B^*\\B&0
  \end{pmatrix}
  =
  \begin{pmatrix}
    0&B^*\\-B&0
  \end{pmatrix}
  =\KB_B.
\end{equation}
\end{proof}

\begin{corollary}[Common magnitude]
\label{cor:common-magnitude}
Both odd branches have the same positive magnitude:
\begin{equation}
  |\DB_B|=|\KB_B|=\DeltaB_B^{1/2}.
\end{equation}
Moreover,
\begin{equation}
  \DB_B^*\DB_B
  =\KB_B^*\KB_B
  =\DeltaB_B.
\end{equation}
\end{corollary}

\begin{proof}
For \(\DB_B\), self-adjointness and \(\DB_B^2=\DeltaB_B\) give the result.
For \(\KB_B\), skew-adjointness and \(\KB_B^2=-\DeltaB_B\) imply
\(\KB_B^*\KB_B=(-\KB_B)\KB_B=\DeltaB_B\).
\end{proof}

\subsection{Balanced equal-carrier channel}

If \(\scrHp\) and \(\scrHm\) are identified with copies of one Hilbert space
\(\mathscr H\), and if \(B=\Lambda=\Lambda^*\), then
\begin{equation}
  \DB_\Lambda=\sigma_x\otimes\Lambda,
  \qquad
  \KB_\Lambda=(i\sigma_y)\otimes\Lambda.
  \label{eq:balanced-branches}
\end{equation}
Thus the Hilbert and Krein branches differ only in the two-dimensional sector
factor. The channel magnitude remains \(|\Lambda|\).
% ===== END INLINED FILE: sections/02_classification.tex =====


% ===== BEGIN INLINED FILE: sections/03_transition_algebra.tex =====
\section{The oriented transition algebra}
\label{sec:algebra}

\subsection{Multiplication table}

\begin{theorem}[Complete \(\Aop/\Vop\) algebra]
\label{thm:AV-algebra}
Let \(B\) satisfy \cref{ass:B}. Then
\begin{align}
  \Aop_B^2&=0,
  &\Vop_B^2&=0,
  \label{eq:AV-nilpotent}\\
  \Aop_B^*&=-\Vop_B,
  &\Vop_B^*&=-\Aop_B,
  \label{eq:AV-star}\\
  \Aop_B^\sharp&=\Vop_B,
  &\Vop_B^\sharp&=\Aop_B.
  \label{eq:AV-ddagger}
\end{align}
Their ordered products are
\begin{equation}
  \Vop_B\Aop_B
  =
  \begin{pmatrix}
    -B^*B&0\\
    0&0
  \end{pmatrix},
  \qquad
  \Aop_B\Vop_B
  =
  \begin{pmatrix}
    0&0\\
    0&-BB^*
  \end{pmatrix}.
  \label{eq:ordered-products}
\end{equation}
On \(\Dom\DeltaB_B\),
\begin{align}
  \{\Aop_B,\Vop_B\}&=-\DeltaB_B,
  \label{eq:anticommutator}\\
  [\Aop_B,\Vop_B]&=\fundJ\DeltaB_B,
  \label{eq:commutator}\\
  \KB_B&=\Aop_B+\Vop_B,
  \label{eq:K-sum}\\
  \DB_B&=\Vop_B-\Aop_B.
  \label{eq:D-difference}
\end{align}
\end{theorem}

\begin{proof}
Nilpotency and the adjoint identities were proved in
\cref{prop:differential-basic}. The ordered products follow by block
multiplication on their natural composition domains. Adding them gives
\(-\DeltaB_B\). Subtracting in the order
\(\Aop_B\Vop_B-\Vop_B\Aop_B\) gives
\(\diag(B^*B,-BB^*)=\fundJ\DeltaB_B\). The final two identities are the
definitions of the two odd branches.
\end{proof}

\subsection{Positive return cycles}

Define the positive ordered return operators
\begin{equation}
  \mathscr C_{+,B}:=-\Vop_B\Aop_B,
  \qquad
  \mathscr C_{-,B}:=-\Aop_B\Vop_B.
  \label{eq:cycles}
\end{equation}
As operators on \(\scrK\),
\begin{equation}
  \mathscr C_{+,B}=
  \begin{pmatrix}
    B^*B&0\\
    0&0
  \end{pmatrix},
  \qquad
  \mathscr C_{-,B}=
  \begin{pmatrix}
    0&0\\
    0&BB^*
  \end{pmatrix}.
  \label{eq:cycles-block}
\end{equation}
Their sector restrictions are the standard partner operators
\begin{equation}
  \restr{\mathscr C_{+,B}}{\scrHp}=B^*B,
  \qquad
  \restr{\mathscr C_{-,B}}{\scrHm}=BB^*.
  \label{eq:partner-restrictions}
\end{equation}
The full Laplacian splits as
\begin{equation}
  \DeltaB_B=\mathscr C_{+,B}+\mathscr C_{-,B}
  \quad\text{on }\Dom\DeltaB_B.
  \label{eq:Delta-cycle-sum}
\end{equation}

\begin{proposition}[Positivity and quadratic forms]
\label{prop:cycle-positivity}
The cycle restrictions \(B^*B\) and \(BB^*\) are positive self-adjoint. For
\(x\in\Dom B\) and \(y\in\Dom B^*\),
\begin{align}
  \|Bx\|^2&=\| (B^*B)^{1/2}x\|^2,
  \label{eq:B-form}\\
  \|B^*y\|^2&=\|(BB^*)^{1/2}y\|^2.
  \label{eq:Bstar-form}
\end{align}
For \(\xi=(x,y)\in\Dom\KB_B=\Dom\DB_B\),
\begin{equation}
  \|\KB_B\xi\|^2
  =\|\DB_B\xi\|^2
  =\|Bx\|^2+\|B^*y\|^2
  =\|\DeltaB_B^{1/2}\xi\|^2.
  \label{eq:transition-energy}
\end{equation}
\end{proposition}

\begin{proof}
Positivity and self-adjointness of \(B^*B\) and \(BB^*\) are standard for a
closed densely defined operator. Their closed quadratic forms are
\(x\mapsto\|Bx\|^2\) and \(y\mapsto\|B^*y\|^2\), giving the first two
identities. The block forms of \(\KB_B\) and \(\DB_B\) give the last identity
directly.
\end{proof}

\subsection{Zero modes}

\begin{proposition}[Common kernel]
\label{prop:kernel}
The kernels of the two odd branches and the partner Laplacian coincide:
\begin{equation}
  \boxed{
  \Ker\DB_B
  =\Ker\KB_B
  =\Ker\DeltaB_B
  =\Ker B\oplus\Ker B^*.}
  \label{eq:kernel}
\end{equation}
The sector kernels are
\begin{equation}
  \Ker(B^*B)=\Ker B,
  \qquad
  \Ker(BB^*)=\Ker B^*.
\end{equation}
\end{proposition}

\begin{proof}
For \(\xi=(x,y)\),
\begin{equation}
  \DB_B\xi=(B^*y,Bx),
  \qquad
  \KB_B\xi=(B^*y,-Bx).
\end{equation}
Either expression vanishes exactly when \(x\in\Ker B\) and
\(y\in\Ker B^*\). The Laplacian claim follows from positivity of the partner
operators.
\end{proof}

\begin{remark}
The two branches have identical zero modes and identical positive magnitude.
Their difference lies entirely in the sector sign: \(\DB_B^2=+\DeltaB_B\)
while \(\KB_B^2=-\DeltaB_B\).
\end{remark}

\subsection{A distinguished reference vector}

Let \(\ket0\in\Dom B\subseteq\scrHp\) be a normalized reference vector,
embedded in \(\scrK\) as \((\ket0,0)\). Define its transition image by
\begin{equation}
  \boxed{\ket\Omega:=\Aop_B\ket0.}
  \label{eq:Omega}
\end{equation}
Explicitly,
\begin{equation}
  \ket\Omega=(0,-B\ket0)\in\scrHm.
  \label{eq:Omega-block}
\end{equation}
No normalization or cyclicity property of \(\ket\Omega\) is implied by this
definition.

\begin{proposition}[One-step image and return]
\label{prop:Omega}
The transition image satisfies
\begin{equation}
  \boxed{\Aop_B\ket\Omega=0.}
  \label{eq:A-Omega}
\end{equation}
If \(\ket0\in\Dom(B^*B)\), then the return transition is defined and obeys
\begin{equation}
  \Vop_B\ket\Omega
  =-B^*B\ket0
  =-\mathscr C_{+,B}\ket0.
  \label{eq:V-Omega}
\end{equation}
Its Hilbert norm and Krein square are
\begin{equation}
  \|\ket\Omega\|^2=\|B\ket0\|^2,
  \qquad
  \kInner{\Omega}{\Omega}=-\|B\ket0\|^2.
  \label{eq:Omega-norms}
\end{equation}
Hence \(\ket\Omega=0\) if and only if \(\ket0\in\Ker B\).
\end{proposition}

\begin{proof}
Equation \eqref{eq:A-Omega} is nilpotency:
\(\Aop_B\ket\Omega=\Aop_B^2\ket0=0\). Under the additional domain
hypothesis,
\(\Vop_B\ket\Omega=\Vop_B\Aop_B\ket0=-B^*B\ket0\). The norm identities
follow from \eqref{eq:Omega-block} and from the negative sign of the Krein
form on \(\scrHm\).
\end{proof}
% ===== END INLINED FILE: sections/03_transition_algebra.tex =====


% ===== BEGIN INLINED FILE: sections/03_symmetry_averaging.tex =====
\section{Compact symmetry averaging and invariant compression}
\label{sec:symmetry-averaging}

The fixed sector projections \(\Pp,\Pm\), a moving physical projection
\(\Pi(x)\), and a symmetry-average projection are different mathematical
objects.  This section records the standard compact-group construction and
fixes the notation \(\Pi_G\) for the latter.

\subsection{The Haar projector}

Let \(G\) be a compact Hausdorff group with normalized Haar measure \(dg\),
and let
\begin{equation}
  \rho:G\longrightarrow\mathcal U(\scrK)
  \label{eq:unitary-representation}
\end{equation}
be a strongly continuous Hilbert-unitary representation.  Define
\begin{equation}
  \Pi_G\xi
  :=\int_G\rho(g)\xi\,dg,
  \qquad \xi\in\scrK,
  \label{eq:Haar-projector}
\end{equation}
where the integral is the Bochner integral of the continuous vector-valued
map \(g\mapsto\rho(g)\xi\).

\begin{theorem}[Compact-group averaging]
\label{thm:Haar-projector}
The operator \(\Pi_G\) is the Hilbert-orthogonal projection onto the invariant
subspace
\begin{equation}
  \scrK^G
  :=\{\xi\in\scrK:\rho(g)\xi=\xi\ \text{for every }g\in G\}.
  \label{eq:invariant-subspace}
\end{equation}
In particular,
\begin{equation}
  \Pi_G^2=\Pi_G,
  \qquad
  \Pi_G^*=\Pi_G,
  \qquad
  \rho(g)\Pi_G=\Pi_G\rho(g)=\Pi_G.
  \label{eq:Haar-properties}
\end{equation}
If in addition
\begin{equation}
  [\rho(g),\fundJ]=0
  \qquad(g\in G),
  \label{eq:J-compatible-representation}
\end{equation}
then
\begin{equation}
  [\Pi_G,\fundJ]=0,
  \label{eq:Haar-J-commute}
\end{equation}
so \(\Pi_G\) reduces the fundamental decomposition and is also
Krein-self-adjoint.
\end{theorem}

\begin{proof}
The estimate
\begin{equation}
  \|\Pi_G\xi\|
  \le\int_G\|\rho(g)\xi\|\,dg
  =\|\xi\|
\end{equation}
shows that \(\Pi_G\) is bounded.  Haar invariance gives, for \(h\in G\),
\begin{equation}
  \rho(h)\Pi_G\xi
  =\int_G\rho(hg)\xi\,dg
  =\Pi_G\xi,
\end{equation}
so \(\Ran\Pi_G\subseteq\scrK^G\).  Conversely, if \(\xi\in\scrK^G\), then
\(\Pi_G\xi=\xi\), hence \(\Ran\Pi_G=\scrK^G\) and \(\Pi_G^2=\Pi_G\).
For \(\xi,\eta\in\scrK\), unitarity and the substitution
\(g\mapsto g^{-1}\) yield
\begin{align}
  \inner{\xi}{\Pi_G\eta}
  &=\int_G\inner{\xi}{\rho(g)\eta}\,dg\\
  &=\int_G\inner{\rho(g^{-1})\xi}{\eta}\,dg
  =\inner{\Pi_G\xi}{\eta},
\end{align}
so \(\Pi_G^*=\Pi_G\).  The remaining intertwining identity follows from
left and right Haar invariance.  Finally, \eqref{eq:J-compatible-representation}
passes through the integral, giving \eqref{eq:Haar-J-commute}; then
\(\Pi_G^\sharp=\fundJ\Pi_G\fundJ=\Pi_G\).
\end{proof}

\begin{corollary}[Finite symmetry group]
\label{cor:finite-group-average}
If \(G\) is finite, then
\begin{equation}
  \Pi_G=\frac1{|G|}\sum_{g\in G}\rho(g).
  \label{eq:finite-group-projector}
\end{equation}
For a finite cubic rotation group, \eqref{eq:finite-group-projector} is the
exact compression onto the trivial cubic representation.
\end{corollary}

\begin{remark}[Distinct projections]
The construction of \(\Pi_G\) does not identify it with the sector projection
\(\Pp\), the active-support projection \(\SuppB\), or a moving projection
\(\Pi(x)\).  Such an identification is additional realization data and must
be checked separately.  In particular, finite cubic invariance and continuum
rotational invariance are different compression problems.
\end{remark}
% ===== END INLINED FILE: sections/03_symmetry_averaging.tex =====


% ===== BEGIN INLINED FILE: sections/03A_reference_state.tex =====
\section{Reference-state geometry and excitation energy}
\label{sec:reference-state}

The transition algebra only requires a normalized reference vector
\(\ket0\in\Dom B\).  This section records an additional result available in
realizations where the same vector is a ground state of a self-adjoint
physical generator.  The statement is a Hilbert-space identity and does not
use the indefinite form.

\subsection{Vacuum fidelity}

Let \(H_+\) be self-adjoint and bounded below on \(\scrHp\).  Assume that
\(\ket0\in\Dom H_+\) is normalized and
\begin{equation}
  H_+\ket0=E_0\ket0,
  \qquad
  E_0=\inf\spec(H_+).
  \label{eq:ground-reference}
\end{equation}
Define
\begin{equation}
  P_0:=\ket0\bra0,
  \qquad
  Q_0:=\idop_+-P_0.
  \label{eq:P0Q0}
\end{equation}
For a normalized \(\psi\in\Dom H_+\), its reference fidelity is
\begin{equation}
  v_\psi:=\inner{\psi}{P_0\psi}
  =|\braket{0}{\psi}|^2\in[0,1].
  \label{eq:reference-fidelity}
\end{equation}
When \(v_\psi<1\), set
\begin{equation}
  \eta_\psi:=\frac{Q_0\psi}{\sqrt{1-v_\psi}}.
  \label{eq:conditional-excitation-vector}
\end{equation}
Then \(\|\eta_\psi\|=1\) and \(\eta_\psi\perp\ket0\).

\begin{theorem}[Reference-fidelity energy decomposition]
\label{thm:reference-fidelity}
For every normalized \(\psi\in\Dom H_+\) with \(v_\psi<1\),
\begin{equation}
  \boxed{
  \inner{\psi}{(H_+-E_0)\psi}
  =(1-v_\psi)
  \inner{\eta_\psi}{(H_+-E_0)\eta_\psi}.}
  \label{eq:reference-fidelity-energy}
\end{equation}
If \(v_\psi=1\), then \(\psi\) lies on the reference ray and the left-hand
side vanishes; in that case the identity is stated separately with right-hand
side zero, without defining \(\eta_\psi\) or \(Z_\psi\).
\end{theorem}

\begin{proof}
Choose \(\theta\in\RR\) such that
\(\braket0\psi=e^{i\theta}\sqrt{v_\psi}\).  Then
\begin{equation}
  \psi=e^{i\theta}\sqrt{v_\psi}\,\ket0
       +\sqrt{1-v_\psi}\,\eta_\psi.
  \label{eq:reference-orthogonal-decomposition}
\end{equation}
Since \((H_+-E_0)\ket0=0\), self-adjointness gives
\begin{equation}
  \inner{\ket0}{(H_+-E_0)\eta_\psi}
  =\inner{(H_+-E_0)\ket0}{\eta_\psi}=0.
\end{equation}
Expanding the quadratic form of \(H_+-E_0\) in
\eqref{eq:reference-orthogonal-decomposition} leaves only the orthogonal
component and proves \eqref{eq:reference-fidelity-energy}.
\end{proof}

\begin{definition}[Conditional excitation scale]
\label{def:conditional-excitation-scale}
For \(v_\psi<1\), define
\begin{equation}
  Z_\psi
  :=\inner{\eta_\psi}{(H_+-E_0)\eta_\psi}\ge0.
  \label{eq:Zpsi}
\end{equation}
Then the excitation energy is
\begin{equation}
  E_{\mathrm{exc}}(\psi)
  :=\inner{\psi}{(H_+-E_0)\psi}
  =Z_\psi(1-v_\psi).
  \label{eq:Eexc-Zv}
\end{equation}
\end{definition}

The scalar \(Z_\psi\) is generally state-dependent.  A state-independent
constant \(Z\) occurs only under additional spectral structure, for example
when \(H_+-E_0\) is constant on the relevant excited subspace.

\subsection{Gap consequence}

\begin{corollary}[Fidelity bound from a spectral gap]
\label{cor:fidelity-gap}
Assume that for some \(m>0\),
\begin{equation}
  H_+-E_0\ge mQ_0
  \label{eq:gap-Q0}
\end{equation}
in the quadratic-form sense.  For \(\psi\in\Dom((H_+-E_0)^{1/2})\), extend the excitation-energy
notation by the closed quadratic form
\begin{equation}
  E_{\mathrm{exc}}(\psi)
  :=\bigl\|(H_+-E_0)^{1/2}\psi\bigr\|^2.
  \label{eq:Eexc-form-domain}
\end{equation}
Then every normalized \(\psi\) in this form domain satisfies
\begin{equation}
  \boxed{
  E_{\mathrm{exc}}(\psi)\ge m(1-v_\psi).}
  \label{eq:fidelity-gap-bound}
\end{equation}
\end{corollary}

\begin{proof}
Apply \eqref{eq:gap-Q0} to \(\psi\) and use
\(\inner\psi{Q_0\psi}=1-v_\psi\).
\end{proof}

If
\begin{equation}
  d_{\mathrm{FS}}([\psi],[\ket0])
  :=\arccos|\braket0\psi|
  \label{eq:FS-distance}
\end{equation}
is the Fubini--Study distance between the ray of \(\psi\) and the reference ray, then
\begin{equation}
  1-v_\psi=\sin^2 d_{\mathrm{FS}}([\psi],[\ket0]),
\end{equation}
and \eqref{eq:fidelity-gap-bound} becomes
\begin{equation}
  E_{\mathrm{exc}}(\psi)
  \ge m\sin^2 d_{\mathrm{FS}}([\psi],[\ket0]).
  \label{eq:FS-gap-bound}
\end{equation}

\subsection{Relation to the transition channel}

The theorem above concerns a self-adjoint physical generator \(H_+\).  It may
be applied to the positive return \(B^*B\), to a Schur--Feshbach effective
operator at a real parameter \(\lambda\in\rho(L_-)\) for which that operator
is self-adjoint, or to another self-adjoint realization, provided the stated
ground-state and domain hypotheses hold.  These cases
must be checked separately.  In particular, the definition
\(\ket\Omega=\Aop_B\ket0\) does not imply that \(\ket0\) is a ground state,
and a ground state in \(\ker B\) has \(\ket\Omega=0\).

Cyclicity is likewise an optional representation property.  If a
\(*\)-algebra \(\mathfrak M\subseteq\mathcal B(\scrHp)\) is specified, the
reference is cyclic when
\begin{equation}
  \overline{\mathfrak M\ket0}=\scrHp.
  \label{eq:cyclic-reference}
\end{equation}
No cyclicity assumption is needed for \cref{thm:reference-fidelity}.
% ===== END INLINED FILE: sections/03A_reference_state.tex =====


% ===== BEGIN INLINED FILE: sections/04_polar_geometry.tex =====
\section{Polar geometry of the transition channel}
\label{sec:polar}

The algebra of the preceding sections separates the orientation of a
transition from its positive magnitude.  The polar decomposition makes that
separation canonical.

\subsection{Polar decomposition and support projections}

Set
\begin{equation}
  |B|:=(B^*B)^{1/2},
  \qquad
  |B^*|:=(BB^*)^{1/2}.
  \label{eq:moduli}
\end{equation}
Both are positive self-adjoint, with
\begin{equation}
  \Dom |B|=\Dom B,
  \qquad
  \Dom |B^*|=\Dom B^*.
\end{equation}

\begin{theorem}[Polar decomposition]
\label{thm:polar-decomposition}
There is a unique partial isometry
\begin{equation}
  u_B:\scrHp\longrightarrow\scrHm
\end{equation}
such that
\begin{equation}
  B=u_B|B|.
  \label{eq:polar-B}
\end{equation}
Its initial and final spaces are
\begin{align}
  \mathscr H_{+,B}^{\mathrm{act}}
  &:=(\Ker B)^\perp
  =\overline{\Ran |B|}
  =\overline{\Ran B^*},
  \label{eq:active-plus}\\
  \mathscr H_{-,B}^{\mathrm{act}}
  &:=(\Ker B^*)^\perp
  =\overline{\Ran |B^*|}
  =\overline{\Ran B}.
  \label{eq:active-minus}
\end{align}
The support projections are
\begin{equation}
  \pi_{+,B}:=u_B^*u_B
  =\idop_+-\operatorname{proj}_{\Ker B},
  \qquad
  \pi_{-,B}:=u_Bu_B^*
  =\idop_--\operatorname{proj}_{\Ker B^*}.
  \label{eq:support-sector}
\end{equation}
Moreover,
\begin{equation}
  |B^*|=u_B|B|u_B^*,
  \qquad
  B^*=|B|u_B^*=u_B^*|B^*|.
  \label{eq:polar-intertwining}
\end{equation}
\end{theorem}

\begin{proof}
This is the polar decomposition theorem for a densely defined closed
operator.  The support identities follow from
\((\Ran B^*)^\perp=\Ker B\) and
\((\Ran B)^\perp=\Ker B^*\).  The remaining identities follow from
functional calculus applied to \(B^*B\) and \(BB^*\); see
\cite{Kato,ReedSimon}.
\end{proof}

Define the total active-support projection
\begin{equation}
  \Pi_B:=
  \begin{pmatrix}
    \pi_{+,B}&0\\
    0&\pi_{-,B}
  \end{pmatrix}.
  \label{eq:PiB}
\end{equation}
Then
\begin{equation}
  \idop_{\scrK}-\Pi_B=\operatorname{proj}_{\Ker\DeltaB_B},
  \qquad
  \Ran\Pi_B=(\Ker\DeltaB_B)^\perp.
  \label{eq:PiB-kernel}
\end{equation}
Thus \(\Pi_B\) removes precisely the kernel data at zero and retains the
part of the two sectors on which the channel has nonzero magnitude.  Any
continuous spectral behavior at zero on the active support remains paired by
\(u_B\).

\subsection{Polar reflection and polar complex structure}

\begin{definition}[Polar reflection and polar complex structure]
\label{def:polar-phases}
The bounded polar reflection and polar complex structure associated with
\(B\) are
\begin{align}
  \mathcal R_B
  &:={}
  \begin{pmatrix}
    0&u_B^*\\
    u_B&0
  \end{pmatrix},
  \label{eq:polar-reflection}\\
  \mathcal I_B
  &:={}
  \begin{pmatrix}
    0&u_B^*\\
    -u_B&0
  \end{pmatrix}.
  \label{eq:polar-complex}
\end{align}
\end{definition}

\begin{theorem}[Supportwise reflection and rotation]
\label{thm:polar-properties}
The polar operators satisfy
\begin{align}
  \mathcal R_B^*&=\mathcal R_B,
  &\mathcal R_B^\sharp&=-\mathcal R_B,
  &\mathcal R_B^2&=\Pi_B,
  \label{eq:R-properties}\\
  \mathcal I_B^*&=-\mathcal I_B,
  &\mathcal I_B^\sharp&=\mathcal I_B,
  &\mathcal I_B^2&=-\Pi_B.
  \label{eq:I-properties}
\end{align}
They vanish on \(\Ker\DeltaB_B\), preserve \(\Ran\Pi_B\), and obey
\begin{equation}
  \boxed{
  \mathcal I_B=\fundJ\mathcal R_B,
  \qquad
  \mathcal R_B=\fundJ\mathcal I_B.}
  \label{eq:IJR}
\end{equation}
On the active support, \(\mathcal R_B\) is a self-adjoint involution and
\(\mathcal I_B\) is a skew-adjoint complex structure.
\end{theorem}

\begin{proof}
The adjoint identities follow directly from the block matrices.  Using
\(u_B^*u_B=\pi_{+,B}\) and \(u_Bu_B^*=\pi_{-,B}\) gives the two square
identities.  Multiplication by \(\fundJ\) changes the sign of the lower-left
block and yields \eqref{eq:IJR}.  Since \(u_B\) and \(u_B^*\) vanish on the
orthogonal complements of their initial spaces, both polar operators vanish
on \(\Ker B\oplus\Ker B^*=\Ker\DeltaB_B\).
\end{proof}

\begin{theorem}[Polar factorization of the odd branches]
\label{thm:polar-factorization}
On the common domain
\begin{equation}
  \Dom\DeltaB_B^{1/2}=\Dom B\oplus\Dom B^*,
\end{equation}
the canonical odd branches factor as
\begin{align}
  \DB_B
  &=\mathcal R_B\DeltaB_B^{1/2}
  =\DeltaB_B^{1/2}\mathcal R_B,
  \label{eq:D-polar}\\
  \KB_B
  &=\mathcal I_B\DeltaB_B^{1/2}
  =\DeltaB_B^{1/2}\mathcal I_B.
  \label{eq:K-polar}
\end{align}
Consequently, \(\mathcal R_B\) and \(\mathcal I_B\) are the bounded polar
phases of \(\DB_B\) and \(\KB_B\), respectively.
\end{theorem}

\begin{proof}
The positive square root is
\begin{equation}
  \DeltaB_B^{1/2}
  =\begin{pmatrix}|B|&0\\0&|B^*|\end{pmatrix}.
\end{equation}
Multiplying by \(\mathcal R_B\) or \(\mathcal I_B\) and using
\eqref{eq:polar-intertwining} gives the stated block matrices.  The same
intertwining identities show that the products agree in either order and
have the claimed common domain.
\end{proof}

\begin{remark}[Generalized inverse formulation]
Let \(\DeltaB_{B,0}^{-1/2}\) be the generalized inverse defined by spectral
calculus from the function \(\lambda\mapsto\lambda^{-1/2}\) on
\((0,\infty)\), with value zero at \(\lambda=0\), and with its natural
spectral domain.  On \(\Ran\DeltaB_B^{1/2}\),
\begin{equation}
  \DB_B\DeltaB_{B,0}^{-1/2}=\restr{\mathcal R_B}{\Ran\DeltaB_B^{1/2}},
  \qquad
  \KB_B\DeltaB_{B,0}^{-1/2}=\restr{\mathcal I_B}{\Ran\DeltaB_B^{1/2}}.
\end{equation}
Because \(\Ran\DeltaB_B^{1/2}\) is dense in the active support
\(\Ran\Pi_B\), these products extend uniquely to bounded operators on that
support.  Extending the active-support operators by zero on
\(\Ker\DeltaB_B\) gives the globally defined polar operators
\(\mathcal R_B\) and \(\mathcal I_B\).
\end{remark}

\subsection{Support-normalized oriented transitions}

\begin{definition}[Normalized transition pair]
\label{def:normalized-AV}
The support-normalized oriented transitions are
\begin{equation}
  \widehat{\Aop}_B
  :=\begin{pmatrix}0&0\\-u_B&0\end{pmatrix},
  \qquad
  \widehat{\Vop}_B
  :=\begin{pmatrix}0&u_B^*\\0&0\end{pmatrix}.
  \label{eq:normalized-AV}
\end{equation}
They are bounded on \(\scrK\).
\end{definition}

\begin{proposition}[Normalized cycle algebra]
\label{prop:normalized-cycle}
The normalized pair satisfies
\begin{align}
  \widehat{\Aop}_B^2&=0,
  &\widehat{\Vop}_B^2&=0,
  &\widehat{\Aop}_B^\sharp&=\widehat{\Vop}_B,
  \label{eq:normalized-basic}\\
  \widehat{\Aop}_B+\widehat{\Vop}_B&=\mathcal I_B,
  &\widehat{\Vop}_B-\widehat{\Aop}_B&=\mathcal R_B,
  \label{eq:normalized-sums}
\end{align}
and
\begin{align}
  -\widehat{\Vop}_B\widehat{\Aop}_B
  &=\begin{pmatrix}\pi_{+,B}&0\\0&0\end{pmatrix},
  \label{eq:normalized-cycle-plus}\\
  -\widehat{\Aop}_B\widehat{\Vop}_B
  &=\begin{pmatrix}0&0\\0&\pi_{-,B}\end{pmatrix}.
  \label{eq:normalized-cycle-minus}
\end{align}
In particular,
\begin{equation}
  -\{\widehat{\Aop}_B,\widehat{\Vop}_B\}=\Pi_B.
  \label{eq:normalized-anticommutator}
\end{equation}
\end{proposition}

\begin{proof}
All statements follow by bounded block multiplication and the support
identities for \(u_B\).
\end{proof}

\subsection{The supportwise rotation group}

\begin{theorem}[Canonical rotation generated by \(\mathcal I_B\)]
\label{thm:rotation-group}
For \(\theta\in\RR\), define
\begin{equation}
  \mathscr E_B(\theta):=e^{\theta\mathcal I_B}.
\end{equation}
Then \(\mathscr E_B(\theta)\) is a strongly continuous Hilbert-unitary group
and
\begin{equation}
  \boxed{
  \mathscr E_B(\theta)
  =(\idop_{\scrK}-\Pi_B)
   +\cos\theta\,\Pi_B
   +\sin\theta\,\mathcal I_B.}
  \label{eq:rotation-exponential}
\end{equation}
On the active support,
\begin{equation}
  \mathscr E_B(\pi/2)=\mathcal I_B,
  \qquad
  \mathscr E_B(\pi)=-\idop,
  \qquad
  \mathscr E_B(2\pi)=\idop.
  \label{eq:quarter-turn}
\end{equation}
Globally,
\begin{equation}
  \mathscr E_B(\pi)=\idop_{\scrK}-2\Pi_B,
\end{equation}
so the zero-mode space is fixed pointwise.
\end{theorem}

\begin{proof}
Since \(\mathcal I_B^*=-\mathcal I_B\), its exponential is Hilbert-unitary.
The identities \(\mathcal I_B^2=-\Pi_B\) and
\(\mathcal I_B(\idop-\Pi_B)=0\) separate the exponential series into the
kernel and active-support parts, giving \eqref{eq:rotation-exponential}.
The special values follow immediately.
\end{proof}

\begin{remark}
The sign \(\mathcal I_B^2=-\Pi_B\) is a normalized two-step return law.  It
is an algebraic cycle sign.  The term \emph{holonomy} will be reserved for a
connection transported around an actual closed path in a parameter space.
\end{remark}
% ===== END INLINED FILE: sections/04_polar_geometry.tex =====


% ===== BEGIN INLINED FILE: sections/05_cohomology_index.tex =====
\section{Cohomology, spectral pairing, and index}
\label{sec:cohomology}

The one-way differential \(d_B\) defines an actual two-term Hilbert complex.
Its cohomology records the zero modes that are not paired by the transition
channel, while the polar partial isometry pairs the entire nonzero spectral
content of the two sectors.

\subsection{The two-term Hilbert complex}

Consider the complex
\begin{equation}
  0\longrightarrow\scrHp
  \xrightarrow{\;B\;}
  \scrHm\longrightarrow0,
  \label{eq:two-term-complex}
\end{equation}
where the differential in degree zero is the closed operator \(B\) with
domain \(\Dom B\).

\begin{definition}[Ordinary and reduced cohomology]
\label{def:cohomology}
The degree-zero and degree-one cohomology spaces are
\begin{equation}
  H_B^0:=\Ker B,
  \qquad
  H_B^1:=\scrHm/\Ran B.
  \label{eq:ordinary-cohomology}
\end{equation}
The reduced degree-one cohomology is
\begin{equation}
  \overline H_B^1:=\scrHm/\overline{\Ran B}.
  \label{eq:reduced-cohomology}
\end{equation}
When \(\Ran B\) is closed, ordinary and reduced degree-one cohomology
coincide.
\end{definition}

\begin{remark}
If \(\Ran B\) is not closed, \(H_B^1\) need not be Hausdorff in the quotient
topology.  The reduced quotient \(\overline H_B^1\) is always a Hilbert
space.  This is the standard distinction for Hilbert complexes
\cite{BruningLesch}.
\end{remark}

\subsection{Hodge decomposition and harmonic representatives}

\begin{theorem}[Two-sector Hodge decomposition]
\label{thm:Hodge-decomposition}
The sector spaces admit the orthogonal decompositions
\begin{align}
  \scrHp
  &=\Ker B\oplus\overline{\Ran B^*},
  \label{eq:Hodge-plus}\\
  \scrHm
  &=\Ker B^*\oplus\overline{\Ran B}.
  \label{eq:Hodge-minus}
\end{align}
Consequently,
\begin{equation}
  \boxed{
  \scrK
  =\Ker\DeltaB_B
   \oplus\overline{\Ran d_B^*}
   \oplus\overline{\Ran d_B}.}
  \label{eq:Hodge-full}
\end{equation}
All sums are Hilbert-orthogonal.  If \(\Ran B\) is closed, the closures in
\eqref{eq:Hodge-full} may be omitted.
\end{theorem}

\begin{proof}
Since \(B\) is densely defined and closed,
\begin{equation}
  (\Ran B^*)^\perp=\Ker B,
  \qquad
  (\Ran B)^\perp=\Ker B^*.
\end{equation}
Taking orthogonal complements gives \eqref{eq:Hodge-plus} and
\eqref{eq:Hodge-minus}.  Since
\begin{equation}
  \Ker\DeltaB_B=\Ker B\oplus\Ker B^*,
\end{equation}
while
\begin{equation}
  \overline{\Ran d_B^*}
  =\overline{\Ran B^*}\oplus\{0\},
  \qquad
  \overline{\Ran d_B}
  =\{0\}\oplus\overline{\Ran B},
\end{equation}
the full decomposition follows.  Closedness of one of \(\Ran B\) and
\(\Ran B^*\) implies closedness of the other.
\end{proof}

\begin{corollary}[Harmonic representation of reduced cohomology]
\label{cor:harmonic-representatives}
There are canonical Hilbert-space isomorphisms
\begin{equation}
  H_B^0\simeq\Ker B,
  \qquad
  \overline H_B^1\simeq\Ker B^*.
  \label{eq:harmonic-cohomology}
\end{equation}
Thus the total reduced cohomology is represented by the harmonic space
\begin{equation}
  H_B^0\oplus\overline H_B^1
  \simeq\Ker\DeltaB_B.
  \label{eq:total-harmonic}
\end{equation}
If \(\Ran B\) is closed, the same statement holds for ordinary cohomology.
\end{corollary}

\begin{proof}
The degree-zero statement is tautological.  In degree one, the orthogonal
projection of \(\scrHm\) onto \(\Ker B^*\) has kernel
\(\overline{\Ran B}\), hence descends to an isometric isomorphism from the
reduced quotient.
\end{proof}

\subsection{Pairing of the nonzero spectra}

Let
\begin{equation}
  H_{+,B}:=B^*B,
  \qquad
  H_{-,B}:=BB^*.
  \label{eq:partner-H}
\end{equation}
The active subspaces from \cref{thm:polar-decomposition} reduce these
operators.  Write
\begin{equation}
  H_{+,B}^{\mathrm{act}}
  :=\restr{H_{+,B}}{\mathscr H_{+,B}^{\mathrm{act}}},
  \qquad
  H_{-,B}^{\mathrm{act}}
  :=\restr{H_{-,B}}{\mathscr H_{-,B}^{\mathrm{act}}}.
\end{equation}

\begin{theorem}[Unitary equivalence of the active partner operators]
\label{thm:partner-unitary-equivalence}
The polar partial isometry restricts to a unitary operator
\begin{equation}
  u_B:\mathscr H_{+,B}^{\mathrm{act}}
  \longrightarrow
  \mathscr H_{-,B}^{\mathrm{act}},
\end{equation}
and
\begin{equation}
  \boxed{
  H_{-,B}^{\mathrm{act}}
  =u_B H_{+,B}^{\mathrm{act}}u_B^*.}
  \label{eq:partner-unitary}
\end{equation}
Equivalently, for every bounded Borel function \(f\),
\begin{equation}
  f(H_{-,B}^{\mathrm{act}})
  =u_B f(H_{+,B}^{\mathrm{act}})u_B^*.
  \label{eq:partner-functional-calculus}
\end{equation}
\end{theorem}

\begin{proof}
On the active spaces, \(u_B^*u_B\) and \(u_Bu_B^*\) are the identity, so
\(u_B\) is unitary.  The polar decomposition gives
\begin{equation}
  BB^*=u_B|B|^2u_B^*=u_B(B^*B)u_B^*
\end{equation}
on the final support, including equality of operator domains.  The Borel
functional-calculus identity follows from unitary equivalence.
\end{proof}

\begin{corollary}[Spectral projection pairing]
\label{cor:spectral-projection-pairing}
Let \(E_+\) and \(E_-\) be the spectral measures of \(B^*B\) and \(BB^*\).
For every Borel set \(S\subset(0,\infty)\),
\begin{equation}
  E_-(S)=u_BE_+(S)u_B^*.
  \label{eq:spectral-projection-pairing}
\end{equation}
Hence
\begin{equation}
  \spec(B^*B)\setminus\{0\}
  =\spec(BB^*)\setminus\{0\},
  \label{eq:nonzero-spectrum}
\end{equation}
and all nonzero eigenvalue multiplicities agree.
\end{corollary}

\begin{corollary}[Explicit eigenvector pairing]
\label{cor:eigenvector-pairing}
If \(\lambda>0\) and
\begin{equation}
  B^*Bx=\lambda x,
\end{equation}
then
\begin{equation}
  y:=\lambda^{-1/2}Bx=u_Bx
\end{equation}
is an eigenvector of \(BB^*\) with eigenvalue \(\lambda\).  Conversely,
\(x=\lambda^{-1/2}B^*y=u_B^*y\).  These maps preserve norms and give an
isomorphism of the corresponding eigenspaces.
\end{corollary}

\begin{proof}
The first claim follows from
\begin{equation}
  BB^*(Bx)=B(B^*Bx)=\lambda Bx.
\end{equation}
Moreover,
\begin{equation}
  \|Bx\|^2=\inner{x}{B^*Bx}=\lambda\|x\|^2.
\end{equation}
The converse is symmetric.
\end{proof}

\subsection{Closed range and a support gap}

\begin{theorem}[Closed-range criterion]
\label{thm:closed-range-criterion}
The following conditions are equivalent:
\begin{enumerate}[label=\textup{(\roman*)}]
  \item \(\Ran B\) is closed;
  \item there exists \(c_B>0\) such that
  \begin{equation}
    \|Bx\|\ge c_B\|x\|
    \quad
    \text{for all }x\in\Dom B\cap(\Ker B)^\perp;
    \label{eq:B-below}
  \end{equation}
  \item on the positive active support, in the quadratic-form sense,
  \begin{equation}
    B^*B\ge c_B^2\pi_{+,B};
    \label{eq:gap-plus}
  \end{equation}
  \item on the negative active support, in the quadratic-form sense,
  \begin{equation}
    BB^*\ge c_B^2\pi_{-,B}.
    \label{eq:gap-minus}
  \end{equation}
\end{enumerate}
When these conditions hold, zero is isolated from the active spectrum of
\(\DeltaB_B\) by at least \(c_B^2\).
\end{theorem}

\begin{proof}
For a closed operator, closedness of the range is equivalent to bounded
invertibility of the restriction
\begin{equation}
  B:\Dom B\cap(\Ker B)^\perp\longrightarrow\Ran B
\end{equation}
when the domain is equipped with its graph norm.  The bounded inverse theorem
gives \eqref{eq:B-below}, and the converse is immediate.  Squaring the norm
inequality yields \eqref{eq:gap-plus}.  The negative-sector statement follows
from \cref{thm:partner-unitary-equivalence}.
\end{proof}

\subsection{Fredholm index and graded zero-mode balance}

\begin{definition}[Fredholm channel]
\label{def:Fredholm-channel}
The closed channel \(B\) is \emph{Fredholm} when
\begin{equation}
  \Ran B\text{ is closed},
  \qquad
  \dim\Ker B<\infty,
  \qquad
  \dim\Ker B^*<\infty.
\end{equation}
Its index is
\begin{equation}
  \ind B:=\dim\Ker B-\dim\Ker B^*.
  \label{eq:index}
\end{equation}
\end{definition}

\begin{theorem}[Cohomological and graded index formulas]
\label{thm:index-formulas}
If \(B\) is Fredholm, then
\begin{equation}
  \boxed{
  \ind B
  =\dim H_B^0-\dim H_B^1
  =\Tr\!\left(\restr{\fundJ}{\Ker\DeltaB_B}\right).}
  \label{eq:index-formulas}
\end{equation}
Equivalently, if \(P_{0,B}:=\operatorname{proj}_{\Ker\DeltaB_B}\), then
\begin{equation}
  \ind B=\Tr(\fundJ P_{0,B}).
  \label{eq:index-projection}
\end{equation}
\end{theorem}

\begin{proof}
Closed range identifies \(H_B^1\) with \(\Ker B^*\) by
\cref{cor:harmonic-representatives}.  On
\(\Ker\DeltaB_B=\Ker B\oplus\Ker B^*\), the grading \(\fundJ\) acts as
\(+1\) on \(\Ker B\) and \(-1\) on \(\Ker B^*\).  Taking the finite
dimensional trace gives the result.
\end{proof}

\begin{corollary}[Fredholm structure of the odd branches]
\label{cor:odd-Fredholm}
If \(B\) is Fredholm, then \(\DB_B\) and \(\KB_B\) have closed range and
finite-dimensional common kernel.  Their ungraded Fredholm indices vanish,
while the grading of their zero-mode space records \(\ind B\) through
\eqref{eq:index-projection}.
\end{corollary}

\subsection{Heat supertrace}

\begin{definition}[Graded trace]
For a block-diagonal trace-class operator
\(X=\diag(X_+,X_-)\), define
\begin{equation}
  \operatorname{Str}_{\fundJ}(X)
  :=\Tr(\fundJ X)
  =\Tr(X_+)-\Tr(X_-).
  \label{eq:supertrace}
\end{equation}
\end{definition}

\begin{theorem}[Heat supertrace index identity]
\label{thm:heat-index}
Suppose \(B\) is Fredholm and both \(e^{-tB^*B}\) and
\(e^{-tBB^*}\) are trace class for every \(t>0\).  Then, for every
\(t>0\),
\begin{equation}
  \boxed{
  \operatorname{Str}_{\fundJ}(e^{-t\DeltaB_B})
  =\Tr(e^{-tB^*B})-\Tr(e^{-tBB^*})
  =\ind B.}
  \label{eq:heat-index}
\end{equation}
\end{theorem}

\begin{proof}
The active restrictions of \(B^*B\) and \(BB^*\) are unitarily equivalent,
so their heat traces agree on the nonzero spectral subspaces.  The remaining
contributions are the zero modes, giving
\begin{equation}
  \dim\Ker B-\dim\Ker B^*=\ind B.
\end{equation}
The value is therefore independent of \(t\).
\end{proof}

\begin{remark}
The heat supertrace isolates the zero-mode imbalance while cancelling the
entire paired nonzero spectrum.  In later realizations, this identity will
separate an index contribution from an ordinary positive heat trace.
\end{remark}
% ===== END INLINED FILE: sections/05_cohomology_index.tex =====


% ===== BEGIN INLINED FILE: sections/06_projector_geometry.tex =====
\section{Moving decompositions and projector curvature}
\label{sec:projector-geometry}

The preceding sections concern one fixed fundamental decomposition.  We now
allow the selected subspace itself to vary over a parameter manifold.  The
resulting tangent operators are automatically grading-odd and therefore fall
under the classification already proved.

\subsection{Norm-\texorpdfstring{\(C^2\)}{C2} projection fields}

Let \(X\) be a finite-dimensional \(C^2\) manifold and let the underlying
positive Hilbert carrier of \(\scrK\) remain fixed.  Suppose
\begin{equation}
  \Pi:X\longrightarrow\mathcal B(\scrK)
  \label{eq:projection-field}
\end{equation}
is twice continuously differentiable in the operator norm and satisfies
\begin{equation}
  \Pi(x)^2=\Pi(x)=\Pi(x)^*,
  \qquad x\in X.
  \label{eq:orthogonal-projection-field}
\end{equation}
Write
\begin{equation}
  \Pi^\perp:=\idop_{\scrK}-\Pi,
  \qquad
  \mathscr E_x:=\Ran\Pi(x),
  \qquad
  \mathscr E_x^\perp:=\Ran\Pi^\perp(x).
  \label{eq:moving-splitting}
\end{equation}
Norm continuity makes the range and corange dimensions locally constant in
the appropriate Hilbert-space sense; hence, on each connected component, the
ranges define a \(C^2\) Hilbert subbundle
\(\mathscr E\subset X\times\scrK\).  At each point, the projection determines
an induced fundamental symmetry
\begin{equation}
  \fundJ_{\Pi}(x):=2\Pi(x)-\idop_{\scrK}
  =\Pi(x)-\Pi^\perp(x).
  \label{eq:moving-J}
\end{equation}
Thus
\begin{equation}
  \scrK=\mathscr E_x\dsum\mathscr E_x^\perp
\end{equation}
with respect to the moving grading \(\fundJ_\Pi(x)\).  The fixed geometry of
\cref{sec:setup} is recovered at any base point for which
\(\Pi(x_0)=\Pp\).

\begin{lemma}[Off-diagonal tangent identity]
\label{lem:projection-tangent}
For every \(C^1\) vector field \(Y\) on \(X\),
\begin{align}
  \Pi(Y\Pi)\Pi&=0,
  &
  \Pi^\perp(Y\Pi)\Pi^\perp&=0,
  \label{eq:tangent-diagonal-zero}
\end{align}
and hence
\begin{equation}
  Y\Pi
  =\Pi^\perp(Y\Pi)\Pi+\Pi(Y\Pi)\Pi^\perp.
  \label{eq:tangent-off-diagonal}
\end{equation}
Moreover, \(Y\Pi\) is Hilbert-self-adjoint.
\end{lemma}

\begin{proof}
Differentiate \(\Pi^2=\Pi\) in the direction \(Y\):
\begin{equation}
  (Y\Pi)\Pi+\Pi(Y\Pi)=Y\Pi.
  \label{eq:differentiate-projection}
\end{equation}
Multiplication by \(\Pi\) on both sides gives
\(\Pi(Y\Pi)\Pi=0\).  Replacing \(\Pi\) by \(\Pi^\perp\) gives the second
identity.  The decomposition \eqref{eq:tangent-off-diagonal} follows by
inserting \(\Pi+\Pi^\perp=\idop\) on both sides of \(Y\Pi\).  Finally,
\((Y\Pi)^*=Y(\Pi^*)=Y\Pi\).
\end{proof}

\begin{definition}[Tangent channel]
\label{def:tangent-channel}
For a \(C^1\) vector field \(Y\), define the bounded channel
\begin{equation}
  B_Y(x):=\Pi^\perp(x)(Y\Pi)(x)\Pi(x)
  :\mathscr E_x\longrightarrow\mathscr E_x^\perp.
  \label{eq:tangent-channel}
\end{equation}
Its adjoint is
\begin{equation}
  B_Y(x)^*=\Pi(x)(Y\Pi)(x)\Pi^\perp(x).
  \label{eq:tangent-channel-adjoint}
\end{equation}
In local coordinates \(x^\mu\), we abbreviate
\begin{equation}
  B_\mu:=B_{\partial_\mu}
  =\Pi^\perp(\partial_\mu\Pi)\Pi.
  \label{eq:Bmu}
\end{equation}
\end{definition}

\begin{theorem}[Tangent realization by the two odd branches]
\label{thm:tangent-realization}
Relative to the moving splitting
\(\scrK=\mathscr E_x\oplus\mathscr E_x^\perp\),
\begin{align}
  \partial_\mu\Pi
  &=
  \begin{pmatrix}
    0&B_\mu^*\\
    B_\mu&0
  \end{pmatrix}
  =\DB_{B_\mu},
  \label{eq:dPi-D}\\
  \fundJ_\Pi\partial_\mu\Pi
  &=
  \begin{pmatrix}
    0&B_\mu^*\\
    -B_\mu&0
  \end{pmatrix}
  =\KB_{B_\mu}.
  \label{eq:JdPi-K}
\end{align}
Equivalently, if
\begin{equation}
  \Aop_\mu:=-d_{B_\mu},
  \qquad
  \Vop_\mu:=d_{B_\mu}^*,
  \label{eq:moving-AV}
\end{equation}
then
\begin{equation}
  \partial_\mu\Pi=\Vop_\mu-\Aop_\mu,
  \qquad
  \fundJ_\Pi\partial_\mu\Pi=\Vop_\mu+\Aop_\mu.
  \label{eq:moving-DK-AV}
\end{equation}
\end{theorem}

\begin{proof}
The first block identity is exactly \cref{lem:projection-tangent} together
with \eqref{eq:tangent-channel-adjoint}.  Multiplication on the left by
\(\fundJ_\Pi=\diag(\idop_{\mathscr E},-\idop_{\mathscr E^\perp})\) changes
the sign of the lower off-diagonal block and gives the second identity.  The
last line follows from the definitions of the two oriented maps.
\end{proof}

\begin{remark}
The two branches now have a direct differential-geometric meaning.  The
Hilbert-self-adjoint operator \(\partial_\mu\Pi\) is the tangent to the
Grassmannian of closed subspaces.  Multiplication by the induced fundamental
symmetry rotates that tangent into the Krein-self-adjoint transition
\(\fundJ_\Pi\partial_\mu\Pi\).
\end{remark}

\subsection{The projected connection}

Let \(C^2(X,\mathscr E)\) denote the \(C^2\) sections of the range bundle.
The trivial ambient bundle \(X\times\scrK\) has the flat derivative \(d\).

\begin{definition}[Grassmann connection]
\label{def:Grassmann-connection}
The projected, or Grassmann, connection on \(\mathscr E\) is
\begin{equation}
  \nabla^\Pi_Y s:=\Pi\,Y s,
  \qquad s\in C^2(X,\mathscr E).
  \label{eq:Grassmann-connection}
\end{equation}
In local coordinates, \(\nabla^\Pi_\mu=\Pi\partial_\mu\).
\end{definition}

\begin{theorem}[Projector curvature]
\label{thm:projector-curvature}
For commuting coordinate vector fields, the curvature of
\(\nabla^\Pi\) is
\begin{align}
  \gaugeF_{\mu\nu}^{\Pi}
  &:=[\nabla^\Pi_\mu,\nabla^\Pi_\nu]
  \label{eq:curvature-def}\\
  &=\Pi[\partial_\mu\Pi,\partial_\nu\Pi]\Pi
  \label{eq:curvature-projector}\\
  &=B_\mu^*B_\nu-B_\nu^*B_\mu
  \quad\text{on }\mathscr E.
  \label{eq:curvature-channel}
\end{align}
Thus the curvature is the antisymmetric part of the mixed closed return
through the complementary bundle.
\end{theorem}

\begin{proof}
Let \(s\in C^2(X,\mathscr E)\), so \(\Pi s=s\).  Then
\begin{align}
  \nabla^\Pi_\mu\nabla^\Pi_\nu s
  &=\Pi\partial_\mu(\Pi\partial_\nu s)\\
  &=\Pi(\partial_\mu\Pi)\partial_\nu s
    +\Pi\partial_\mu\partial_\nu s.
\end{align}
Differentiating \(s=\Pi s\) gives
\(\partial_\nu s=(\partial_\nu\Pi)s+\Pi\partial_\nu s\).  The term with
\(\Pi\partial_\nu s\) vanishes because
\(\Pi(\partial_\mu\Pi)\Pi=0\).  Antisymmetrizing in \(\mu,\nu\) therefore
gives \eqref{eq:curvature-projector}.  In the block form of
\cref{thm:tangent-realization}, the positive-sector block of
\([\partial_\mu\Pi,\partial_\nu\Pi]\) is
\(B_\mu^*B_\nu-B_\nu^*B_\mu\).
\end{proof}

\begin{definition}[Mixed return operators]
\label{def:mixed-return}
For two tangent directions, define the positive-sector mixed return
\begin{equation}
  \mathscr C_{\mu\nu}^{\Pi}
  :=-\restr{\Vop_\mu\Aop_\nu}{\mathscr E}
  =B_\mu^*B_\nu.
  \label{eq:mixed-return}
\end{equation}
\end{definition}

\begin{corollary}[Curvature as an antisymmetrized closed excursion]
\label{cor:curvature-AV}
The curvature may be written
\begin{align}
  \gaugeF_{\mu\nu}^{\Pi}
  &=\mathscr C_{\mu\nu}^{\Pi}-\mathscr C_{\nu\mu}^{\Pi}
  \label{eq:curvature-mixed-return}\\
  &=\restr{\Vop_\nu\Aop_\mu-\Vop_\mu\Aop_\nu}{\mathscr E}.
  \label{eq:curvature-AV}
\end{align}
\end{corollary}

\begin{remark}
For a single fixed channel, the ordered return
\(-\Vop_B\Aop_B=B^*B\) is positive.  Curvature is not one such return but
the antisymmetric difference of two mixed returns associated with distinct
tangent directions.  Nonzero curvature therefore measures the failure of
these directional closed excursions to agree under interchange.
\end{remark}

\subsection{Local frames and gauge covariance}

Suppose that \(\mathscr E\) has finite rank \(r\) over a coordinate
neighborhood \(U\subseteq X\).  Choose a \(C^2\) orthonormal frame written as
an isometry
\begin{equation}
  W:U\longrightarrow\mathcal B(\CC^r,\scrK),
  \qquad
  W^*W=\idop_r,
  \qquad
  WW^*=\Pi.
  \label{eq:local-frame}
\end{equation}
We use the anti-Hermitian convention for the local connection matrix.

\begin{definition}[Local gauge data]
\label{def:local-gauge-data}
The connection and curvature in the frame \(W\) are
\begin{align}
  \gaugeA_\mu^{(W)}&:=W^*\partial_\mu W,
  \qquad
  (\gaugeA_\mu^{(W)})^*=-\gaugeA_\mu^{(W)},
  \label{eq:local-connection}\\
  \gaugeF_{\mu\nu}^{(W)}
  &:=\partial_\mu\gaugeA_\nu^{(W)}
    -\partial_\nu\gaugeA_\mu^{(W)}
    +[\gaugeA_\mu^{(W)},\gaugeA_\nu^{(W)}].
  \label{eq:local-curvature}
\end{align}
\end{definition}

\begin{theorem}[Frame covariance]
\label{thm:frame-covariance}
The local curvature is the frame representation of the projector curvature:
\begin{equation}
  \boxed{
  \gaugeF_{\mu\nu}^{(W)}
  =W^*\gaugeF_{\mu\nu}^{\Pi}W
  =W^*(B_\mu^*B_\nu-B_\nu^*B_\mu)W.}
  \label{eq:frame-curvature}
\end{equation}
If \(W'=Wg\) for a \(C^2\) map \(g:U\to\mathrm U(r)\), then
\begin{align}
  \gaugeA_\mu^{(W')}
  &=g^{-1}\gaugeA_\mu^{(W)}g+g^{-1}\partial_\mu g,
  \label{eq:gauge-transform-A}\\
  \gaugeF_{\mu\nu}^{(W')}
  &=g^{-1}\gaugeF_{\mu\nu}^{(W)}g.
  \label{eq:gauge-transform-F}
\end{align}
\end{theorem}

\begin{proof}
The transformation law for \(\gaugeA\) follows immediately from
\((Wg)^*\partial_\mu(Wg)\) and \(g^*=g^{-1}\).  Substitution into
\eqref{eq:local-curvature} gives the curvature covariance.  For the projector
formula, differentiate \(\Pi=WW^*\), use \(W^*W=\idop_r\), and simplify the
standard identity
\begin{equation}
  d(W^*dW)+(W^*dW)\wedge(W^*dW)
  =W^*(d\Pi\wedge d\Pi)W.
\end{equation}
Its coordinate form is exactly \eqref{eq:frame-curvature}.
\end{proof}

\begin{definition}[Holonomy]
\label{def:holonomy}
Let \(\gamma:[0,1]\to U\) be a piecewise \(C^1\) loop.  Parallel transport in
the frame \(W\) is the solution of
\begin{equation}
  \dot c(t)+\gaugeA^{(W)}_{\gamma(t)}(\dot\gamma(t))c(t)=0.
  \label{eq:parallel-transport}
\end{equation}
The resulting unitary endomorphism of the initial fiber is the holonomy
\begin{equation}
  \operatorname{Hol}_\gamma(\nabla^\Pi)\in\mathrm U(r).
  \label{eq:holonomy}
\end{equation}
\end{definition}

\begin{remark}
This is the point at which the term \emph{holonomy} acquires its standard
geometric meaning.  The fixed-channel identity \(\mathcal I_B^2=-\Pi_B\) of
\cref{sec:polar} is a supportwise cycle sign; holonomy requires a connection
and transport around an actual loop in the base manifold.
\end{remark}

\subsection{Fiberwise inheritance of the fixed-channel algebra}

At every \(x\in X\) and every tangent direction \(Y\in T_xX\), the bounded
channel
\begin{equation}
  B_Y(x):\mathscr E_x\longrightarrow\mathscr E_x^\perp
\end{equation}
inherits all fixed-channel identities of \cref{sec:classification,sec:algebra,sec:polar}.  In particular,
\begin{equation}
  \Delta_{B_Y}
  =\begin{pmatrix}B_Y^*B_Y&0\\0&B_YB_Y^*\end{pmatrix},
  \qquad
  (\fundJ_\Pi Y\Pi)^2=-\Delta_{B_Y}.
  \label{eq:fiberwise-Laplacian}
\end{equation}
The moving-projection layer therefore adds no new local odd algebra.  It
selects a geometrically canonical family of channels and couples distinct
directions through their mixed returns.

\begin{equation}
  \boxed{
  \text{projection tangent}
  \xrightarrow{\ \fundJ_\Pi\ }
  \text{Krein transition},
  \qquad
  \text{antisymmetrized mixed return}
  =\text{curvature}.}
  \label{eq:projector-summary}
\end{equation}
% ===== END INLINED FILE: sections/06_projector_geometry.tex =====


% ===== BEGIN INLINED FILE: sections/06A_characteristic_classes.tex =====
\section{Characteristic classes of the projector bundle}
\label{sec:characteristic-classes}

The projector curvature of \cref{sec:projector-geometry} has the standard
Chern--Weil consequences when the selected bundle has finite rank.  This
section records those consequences and expresses them in terms of the mixed
transition returns.

Let \(X\) be a smooth manifold and let
\(\Pi:X\to\mathcal B(\scrK)\) be a smooth finite-rank orthogonal projection
field.  The range bundle is \(\mathscr E=\Ran\Pi\to X\), equipped with the
Grassmann connection \(\nabla^\Pi=\Pi d\).  Its anti-Hermitian curvature is
\begin{equation}
  \gaugeF^\Pi
  =\Pi(d\Pi\wedge d\Pi)\Pi.
  \label{eq:curvature-form}
\end{equation}
In coordinates this is the operator-valued two-form whose components are
\begin{equation}
  \gaugeF_{\mu\nu}^{\Pi}
  =B_\mu^*B_\nu-B_\nu^*B_\mu.
  \label{eq:curvature-return-components}
\end{equation}

\subsection{Chern--Weil forms}

\begin{theorem}[Chern character from closed returns]
\label{thm:chern-character}
The differential form
\begin{equation}
  \operatorname{ch}(\mathscr E,\nabla^\Pi)
  :=\Tr_{\mathscr E}
  \exp\!\left(\frac{i}{2\pi}\gaugeF^\Pi\right)
  \label{eq:chern-character}
\end{equation}
is closed.  Its de Rham cohomology class is independent of the chosen
connection and equals the Chern character of \(\mathscr E\).  In particular,
the first Chern form is
\begin{equation}
  c_1(\mathscr E,\nabla^\Pi)
  =\frac{i}{2\pi}\Tr_{\mathscr E}(\gaugeF^\Pi)
  =\frac{i}{2\pi}\Tr_{\mathscr E}
    \bigl(\Pi(d\Pi\wedge d\Pi)\Pi\bigr).
  \label{eq:first-chern-form}
\end{equation}
In local coordinates,
\begin{equation}
  c_1(\mathscr E,\nabla^\Pi)_{\mu\nu}
  =\frac{i}{2\pi}
  \Tr_{\mathscr E}
  \bigl(B_\mu^*B_\nu-B_\nu^*B_\mu\bigr).
  \label{eq:first-chern-mixed-return}
\end{equation}
\end{theorem}

\begin{proof}
The Grassmann connection is a unitary connection on the finite-rank Hermitian
bundle \(\mathscr E\).  The Chern--Weil theorem applies to the invariant
polynomials \(A\mapsto\Tr(A^k)\), proving closedness and connection
independence of the resulting cohomology classes
\cite{KobayashiNomizu,Husemoller}.  Formula
\eqref{eq:first-chern-form} is the degree-two component of
\eqref{eq:chern-character}.  Substitution of
\eqref{eq:curvature-return-components} gives
\eqref{eq:first-chern-mixed-return}.
\end{proof}

Thus the same mixed returns that generate local curvature also generate the
characteristic forms of the selected bundle.  Higher-degree terms are
obtained from
\begin{equation}
  \operatorname{ch}_k(\mathscr E,\nabla^\Pi)
  =\frac{1}{k!}
  \left(\frac{i}{2\pi}\right)^k
  \Tr_{\mathscr E}\bigl((\gaugeF^\Pi)^{\wedge k}\bigr).
  \label{eq:higher-chern-character}
\end{equation}

\subsection{Integral flux in rank one}

\begin{corollary}[Line-bundle flux quantization]
\label{cor:line-flux}
Suppose \(\mathscr E\to X\) has rank one and let \(\Sigma\) be a closed
oriented smooth two-cycle.  Then
\begin{equation}
  \boxed{
  \frac{i}{2\pi}\int_\Sigma \gaugeF^\Pi\in\ZZ.}
  \label{eq:line-flux-integral}
\end{equation}
The integer is the pairing
\(\langle c_1(\mathscr E),[\Sigma]\rangle\).
\end{corollary}

\begin{proof}
For a Hermitian line bundle, the de Rham class of
\((i/2\pi)\gaugeF^\Pi\) is the image of the integral class
\(c_1(\mathscr E)\in H^2(X;\ZZ)\).  Pairing with the integral homology class
\([\Sigma]\) gives an integer.
\end{proof}

The connection convention in this paper is anti-Hermitian.  If the
curvature is converted to the Hermitian two-form
\begin{equation}
  \gaugeF_{\mathrm H}:=i\gaugeF,
\end{equation}
then \eqref{eq:line-flux-integral} becomes
\begin{equation}
  \frac{1}{2\pi}\int_\Sigma\gaugeF_{\mathrm H}\in\ZZ.
  \label{eq:physics-flux-integral}
\end{equation}

\begin{remark}[Holonomy versus characteristic number]
The holonomy of a connection around a closed loop need not be an integral
multiple of \(2\pi\); even a contractible loop can have nontrivial holonomy
when curvature is present.  The integral statement is instead the pairing of
the curvature class with a closed two-cycle, as in
\eqref{eq:line-flux-integral}.
\end{remark}

\subsection{Transition-algebra form}

Using the mixed-return notation of \cref{def:mixed-return},
\begin{equation}
  \gaugeF_{\mu\nu}^{\Pi}
  =\mathscr C_{\mu\nu}^{\Pi}
   -\mathscr C_{\nu\mu}^{\Pi}.
\end{equation}
Consequently, every Chern--Weil density is a trace polynomial in
antisymmetrized closed excursions.  For example,
\begin{equation}
  \operatorname{ch}_2(\mathscr E,\nabla^\Pi)
  =\frac{1}{2}
   \left(\frac{i}{2\pi}\right)^2
   \Tr_{\mathscr E}(\gaugeF^\Pi\wedge\gaugeF^\Pi).
  \label{eq:ch2-return}
\end{equation}
This is an exact bundle-theoretic statement.  Identifying a particular
characteristic number with a physical charge requires additional realization
data and is not part of the projector theorem itself.
% ===== END INLINED FILE: sections/06A_characteristic_classes.tex =====


% ===== BEGIN INLINED FILE: sections/07_effective_dynamics.tex =====
\section{Schur--Feshbach reduction and effective dynamics}
\label{sec:effective-dynamics}

The preceding curvature is built from undressed mixed returns
\(B_\mu^*B_\nu\).  We now eliminate one sector from a block generator.  The
same transition channel then reappears as a resolvent-dressed closed return.
To keep the operator domains fully controlled, this pass states the exact
bounded-coupling theorem.  Form-coupled and relatively bounded extensions
require additional hypotheses and are not assumed here.

\subsection{Two sign conventions for one block channel}

\begin{assumption}[Diagonal generators and bounded coupling]
\label{ass:block-generator}
Let
\begin{equation}
  L_+:\Dom L_+\subseteq\scrHp\to\scrHp,
  \qquad
  L_-:\Dom L_-\subseteq\scrHm\to\scrHm
  \label{eq:diagonal-generators}
\end{equation}
be self-adjoint, and let
\begin{equation}
  B\in\mathcal B(\scrHp,\scrHm).
  \label{eq:bounded-coupling}
\end{equation}
For \(\varepsilon\in\{+1,-1\}\), define
\begin{equation}
  \mathbb L^{(\varepsilon)}
  :=
  \begin{pmatrix}
    L_+&B^*\\
    \varepsilon B&L_-
  \end{pmatrix},
  \qquad
  \Dom\mathbb L^{(\varepsilon)}
  :=\Dom L_+\oplus\Dom L_-.
  \label{eq:block-generator}
\end{equation}
\end{assumption}

\begin{proposition}[Hilbert and Krein block branches]
\label{prop:block-branches}
Let
\begin{equation}
  \mathbb L_0:=\diag(L_+,L_-).
\end{equation}
Then
\begin{align}
  \mathbb L^{(+1)}
  &=\mathbb L_0+\DB_B,
  &
  (\mathbb L^{(+1)})^*&=\mathbb L^{(+1)},
  \label{eq:block-Hilbert-branch}\\
  \mathbb L^{(-1)}
  &=\mathbb L_0+\KB_B,
  &
  (\mathbb L^{(-1)})^\sharp&=\mathbb L^{(-1)}.
  \label{eq:block-Krein-branch}
\end{align}
Thus the same channel \(B\) produces a Hilbert-self-adjoint coupling for
\(\varepsilon=+1\) and a Krein-self-adjoint coupling for
\(\varepsilon=-1\).
\end{proposition}

\begin{proof}
Because \(B\) is bounded, both block matrices are bounded perturbations of
\(\mathbb L_0\) on its natural domain.  The off-diagonal part for
\(\varepsilon=+1\) is \(\DB_B\), which is Hilbert-self-adjoint.  The
off-diagonal part for \(\varepsilon=-1\) is \(\KB_B\), which is
Krein-self-adjoint.  The diagonal operator commutes with \(\fundJ\) and is
self-adjoint in both sectors, so it is also Krein-self-adjoint.
\end{proof}

\subsection{The operator-valued Schur complement}

Let \(z\in\rho(L_-)\) and set
\begin{equation}
  R_-(z):=(L_--z)^{-1}.
  \label{eq:Rminus}
\end{equation}
The complementary-sector return is the bounded operator
\begin{equation}
  \Sigma_B(z):=B^*R_-(z)B
  \in\mathcal B(\scrHp).
  \label{eq:self-energy}
\end{equation}

\begin{definition}[Effective operator and Schur pencil]
\label{def:effective-operator}
For \(\varepsilon\in\{+1,-1\}\), define
\begin{align}
  L_{\mathrm{eff}}^{(\varepsilon)}(z)
  &:=L_+-\varepsilon\Sigma_B(z),
  \label{eq:effective-operator}\\
  S_\varepsilon(z)
  &:=L_{\mathrm{eff}}^{(\varepsilon)}(z)-z
  =L_+-z-\varepsilon B^*R_-(z)B
  \label{eq:Schur-pencil}
\end{align}
on \(\Dom L_+\).
\end{definition}

\begin{theorem}[Schur factorization and compressed resolvent]
\label{thm:Schur-factorization}
For \(z\in\rho(L_-)\),
\begin{equation}
  \mathbb L^{(\varepsilon)}-z
  =
  \begin{pmatrix}
    \idop_+&B^*R_-(z)\\
    0&\idop_-
  \end{pmatrix}
  \begin{pmatrix}
    S_\varepsilon(z)&0\\
    \varepsilon B&L_--z
  \end{pmatrix}.
  \label{eq:Schur-factorization}
\end{equation}
Consequently,
\begin{equation}
  z\in\rho(\mathbb L^{(\varepsilon)})
  \quad\Longleftrightarrow\quad
  0\in\rho(S_\varepsilon(z)).
  \label{eq:Schur-invertibility}
\end{equation}
When these equivalent conditions hold,
\begin{equation}
  (\mathbb L^{(\varepsilon)}-z)^{-1}
  =
  \begin{pmatrix}
    S_\varepsilon(z)^{-1}
    &-S_\varepsilon(z)^{-1}B^*R_-(z)\\[2mm]
    -\varepsilon R_-(z)B S_\varepsilon(z)^{-1}
    &R_-(z)+\varepsilon R_-(z)B S_\varepsilon(z)^{-1}B^*R_-(z)
  \end{pmatrix}.
  \label{eq:block-resolvent}
\end{equation}
In particular,
\begin{equation}
  \boxed{
  \Pp(\mathbb L^{(\varepsilon)}-z)^{-1}\Pp
  =\bigl(L_{\mathrm{eff}}^{(\varepsilon)}(z)-z\bigr)^{-1}.}
  \label{eq:compressed-resolvent}
\end{equation}
\end{theorem}

\begin{proof}
Multiplication of the two block factors gives the upper-left entry
\begin{equation}
  S_\varepsilon(z)+\varepsilon B^*R_-(z)B=L_+-z,
\end{equation}
the upper-right entry \(B^*R_-(z)(L_--z)=B^*\), and the stated lower row.
The first factor is boundedly invertible.  Since \(L_--z\) is invertible, the
second factor is invertible exactly when \(S_\varepsilon(z)\) is invertible.  Direct
block multiplication then verifies \eqref{eq:block-resolvent}; its
positive-sector corner gives \eqref{eq:compressed-resolvent}.
\end{proof}

\subsection{The self-energy as a dressed closed excursion}

Embed the negative resolvent into \(\scrK\) by
\begin{equation}
  \widetilde R_-(z):=
  \begin{pmatrix}
    0&0\\
    0&R_-(z)
  \end{pmatrix}.
  \label{eq:embedded-resolvent}
\end{equation}

\begin{proposition}[Oriented excursion formula]
\label{prop:oriented-self-energy}
The self-energy is the return obtained by leaving the positive sector through
\(\Aop_B\), propagating in the complementary sector, and returning through
\(\Vop_B\):
\begin{equation}
  \boxed{
  \Sigma_B(z)
  =-\restr{\Vop_B\widetilde R_-(z)\Aop_B}{\scrHp}.}
  \label{eq:self-energy-AV}
\end{equation}
Hence
\begin{equation}
  L_{\mathrm{eff}}^{(\varepsilon)}(z)
  =L_++\varepsilon
  \restr{\Vop_B\widetilde R_-(z)\Aop_B}{\scrHp}.
  \label{eq:effective-AV}
\end{equation}
\end{proposition}

\begin{proof}
For \(x\in\scrHp\),
\begin{equation}
  \Aop_Bx=(0,-Bx),
  \qquad
  \widetilde R_-(z)\Aop_Bx=(0,-R_-(z)Bx),
\end{equation}
and applying \(\Vop_B\) gives
\((-B^*R_-(z)Bx,0)\).  Restriction to the positive sector and multiplication
by \(-1\) yields \(\Sigma_B(z)\).
\end{proof}

\begin{remark}
The sign in the effective operator is not a matter of notation.  It is fixed
by the operator category of the full block generator:
\begin{align}
  L_{\mathrm{eff}}^{(+1)}(z)
  &=L_+-\Sigma_B(z),
  \label{eq:Leff-plus}\\
  L_{\mathrm{eff}}^{(-1)}(z)
  &=L_++\Sigma_B(z).
  \label{eq:Leff-minus}
\end{align}
The Hilbert/Hodge and Krein/rotational branches therefore induce opposite
self-energy signs after the same complementary sector is eliminated.
\end{remark}

\subsection{Analytic and positivity properties of the return}

\begin{theorem}[Nevanlinna property of the self-energy]
\label{thm:self-energy-Herglotz}
The map
\begin{equation}
  \Sigma_B:\CC\setminus\RR\longrightarrow\mathcal B(\scrHp)
\end{equation}
is analytic and satisfies
\begin{equation}
  \Sigma_B(\overline z)=\Sigma_B(z)^*.
  \label{eq:self-energy-symmetry}
\end{equation}
For \(z\in\CC_+:=\{z:\operatorname{Im}z>0\}\),
\begin{align}
  \operatorname{Im}\Sigma_B(z)
  &:=\frac{\Sigma_B(z)-\Sigma_B(z)^*}{2i}
  \label{eq:Im-Sigma-def}\\
  &=(\operatorname{Im}z)
    B^*R_-(z)R_-(z)^*B
  \ge0.
  \label{eq:Im-Sigma}
\end{align}
Thus \(\Sigma_B\) is an operator-valued Herglotz--Nevanlinna function.
\end{theorem}

\begin{proof}
Analyticity and the adjoint symmetry follow from the resolvent of the
self-adjoint operator \(L_-\).  The resolvent identity gives
\begin{equation}
  R_-(z)-R_-(\overline z)
  =(z-\overline z)R_-(z)R_-(\overline z).
\end{equation}
Since \(R_-(\overline z)=R_-(z)^*\), multiplication by \(B^*\) and \(B\)
gives \eqref{eq:Im-Sigma}.  Positivity is immediate.
\end{proof}

\begin{corollary}[Opposite dissipative signs]
\label{cor:effective-imaginary-sign}
For \(z\in\CC_+\),
\begin{align}
  \operatorname{Im}L_{\mathrm{eff}}^{(+1)}(z)
  &=-\operatorname{Im}\Sigma_B(z)\le0,
  \label{eq:Leff-plus-imag}\\
  \operatorname{Im}L_{\mathrm{eff}}^{(-1)}(z)
  &=+\operatorname{Im}\Sigma_B(z)\ge0.
  \label{eq:Leff-minus-imag}
\end{align}
\end{corollary}

\begin{corollary}[Real energies below the complementary spectrum]
\label{cor:static-self-energy}
If \(\lambda\in\RR\) satisfies
\(\lambda<\inf\spec L_-\), then
\begin{equation}
  R_-(\lambda)>0,
  \qquad
  \Sigma_B(\lambda)\ge0.
  \label{eq:static-positive-self-energy}
\end{equation}
Accordingly, the Hilbert branch lowers the effective quadratic form by
\(\Sigma_B(\lambda)\), whereas the Krein branch raises it by the same return.
\end{corollary}

\subsection{One mechanism in two geometric forms}

The variable-projection curvature of \cref{sec:projector-geometry} and the
Schur self-energy of the present section are built from the same algebraic
primitive.

\begin{theorem}[Undressed and dressed closed returns]
\label{thm:two-return-forms}
Let \(B_\mu\) be the tangent channels of a norm-\(C^2\) projection field.  Then
\begin{equation}
  \gaugeF_{\mu\nu}^{\Pi}
  =B_\mu^*B_\nu-B_\nu^*B_\mu
  \label{eq:undressed-return}
\end{equation}
is the antisymmetrized undressed mixed return.  For a block generator with
complementary dynamics \(L_-\),
\begin{equation}
  \Sigma_B(z)=B^*(L_--z)^{-1}B
  \label{eq:dressed-return}
\end{equation}
is the resolvent-dressed return through that sector.  Both are closed
positive-sector excursions of the form
\begin{equation}
  \mathscr H_+
  \xrightarrow{\ B\ }
  \mathscr H_-
  \xrightarrow{\ \text{intermediate operation}\ }
  \mathscr H_-
  \xrightarrow{\ B^*\ }
  \mathscr H_+.
  \label{eq:excursion-template}
\end{equation}
\end{theorem}

\begin{proof}
The first identity is \cref{thm:projector-curvature}; the second is the
definition of the Schur self-energy.  Their common source and target follow
from the channel types.
\end{proof}

\begin{equation}
  \boxed{
  \begin{aligned}
  \text{antisymmetrized mixed return}
  &\longrightarrow \text{connection curvature},\\
  \text{resolvent-dressed return}
  &\longrightarrow \text{effective dynamics}.
  \end{aligned}}
  \label{eq:pass3-summary}
\end{equation}
% ===== END INLINED FILE: sections/07_effective_dynamics.tex =====


% ===== BEGIN INLINED FILE: sections/08_boundary_response.tex =====
\section{Boundary response and passive realizations}
\label{sec:boundary-response}

The Schur--Feshbach reduction of \cref{sec:effective-dynamics} identifies the
operator seen after the complementary sector has been eliminated.  We now
separate that reduction from a second issue: how a self-adjoint realization
is coupled to a finite or otherwise controlled port space.  This distinction
is essential.  Positivity of a response follows from Hilbert-space
self-adjointness and does not follow from Krein self-adjointness alone.
The operator-valued Nevanlinna, Schur, and reproducing-kernel constructions
used in this section are standard; see
\cite{Donoghue,Aronszajn,deBrangesRovnyak}.

\subsection{Compressed resolvents}

Let \(\scrX\) and \(\scrE\) be complex Hilbert spaces.  The space
\(\scrE\) will be called the port space; no finite-dimensionality assumption
is needed unless stated explicitly.

\begin{assumption}[Self-adjoint response realization]
\label{ass:selfadjoint-response}
Let
\begin{equation}
  H:\Dom H\subseteq\scrX\longrightarrow\scrX
  \label{eq:response-H}
\end{equation}
be self-adjoint and let
\begin{equation}
  \Gamma\in\mathcal B(\scrE,\scrX).
  \label{eq:response-Gamma}
\end{equation}
For \(z\in\CC\setminus\RR\), write
\begin{equation}
  R_H(z):=(H-z)^{-1}.
  \label{eq:response-resolvent}
\end{equation}
\end{assumption}

\begin{definition}[Boundary response]
\label{def:boundary-response}
The compressed resolvent, or boundary response, of the realization
\((H,\Gamma)\) is
\begin{equation}
  \boxed{
  \Mresp(z):=\Gamma^*R_H(z)\Gamma,
  \qquad z\in\CC\setminus\RR.}
  \label{eq:boundary-response}
\end{equation}
\end{definition}

The word ``boundary'' refers to the role of \(\Gamma\) as an input or trace
channel.  The definition itself is abstract and applies equally to interior,
finite-rank, and regularized distributional inputs.

\begin{theorem}[Herglotz--Nevanlinna response]
\label{thm:response-Herglotz}
The response \(\Mresp\) is analytic on \(\CC\setminus\RR\), satisfies
\begin{equation}
  \Mresp(\overline z)=\Mresp(z)^*,
  \label{eq:response-symmetry}
\end{equation}
and, for \(z\in\Cplus\),
\begin{equation}
  \boxed{
  \operatorname{Im}\Mresp(z)
  =(\operatorname{Im}z)\,
  \Gamma^*R_H(z)R_H(z)^*\Gamma
  \geq0.}
  \label{eq:response-positive-imaginary}
\end{equation}
Thus \(\Mresp\) is an operator-valued Herglotz--Nevanlinna function.
\end{theorem}

\begin{proof}
Analyticity and the adjoint symmetry are inherited from the resolvent of the
self-adjoint operator \(H\).  The resolvent identity gives
\begin{equation}
  R_H(z)-R_H(\overline z)
  =(z-\overline z)R_H(z)R_H(\overline z).
\end{equation}
Since \(R_H(\overline z)=R_H(z)^*\), compression by \(\Gamma\) gives
\eqref{eq:response-positive-imaginary}.
\end{proof}

\begin{definition}[Nevanlinna kernel]
\label{def:Nevanlinna-kernel}
For \(z,w\in\Cplus\), define
\begin{equation}
  \KM(z,w)
  :=\frac{\Mresp(z)-\Mresp(w)^*}{z-\overline w}.
  \label{eq:Nevanlinna-kernel}
\end{equation}
\end{definition}

\begin{theorem}[Positive response kernel]
\label{thm:positive-response-kernel}
The kernel \(\KM\) is positive in the sense that, for every finite collection
\(z_1,\ldots,z_N\in\Cplus\) and \(h_1,\ldots,h_N\in\scrE\),
\begin{equation}
  \sum_{i,j=1}^{N}
  \inner{h_i}{\KM(z_i,z_j)h_j}_{\scrE}\geq0.
  \label{eq:Nevanlinna-kernel-positive}
\end{equation}
More precisely,
\begin{equation}
  \boxed{
  \KM(z,w)=\Gamma^*R_H(z)R_H(w)^*\Gamma.}
  \label{eq:Nevanlinna-kernel-factor}
\end{equation}
\end{theorem}

\begin{proof}
The two-point resolvent identity gives
\begin{equation}
  R_H(z)-R_H(\overline w)
  =(z-\overline w)R_H(z)R_H(\overline w).
\end{equation}
After compression, this is \eqref{eq:Nevanlinna-kernel-factor}.  Hence the
left-hand side of \eqref{eq:Nevanlinna-kernel-positive} equals
\begin{equation}
  \left\|
  \sum_{j=1}^{N}R_H(z_j)^*\Gamma h_j
  \right\|_{\scrX}^{2}.
\end{equation}
\end{proof}

\begin{corollary}[Canonical transition response]
\label{cor:canonical-transition-response}
For the transition channel \(B\) of \cref{sec:setup}, the positive return
operator
\begin{equation}
  H_{+,B}:=B^*B
  \label{eq:positive-return-H}
\end{equation}
is self-adjoint and nonnegative.  Therefore every bounded input
\(\Gamma:\scrE\to\scrHp\) defines the canonical response
\begin{equation}
  M_{B,\Gamma}(z)
  :=\Gamma^*(B^*B-z)^{-1}\Gamma,
  \label{eq:canonical-transition-response}
\end{equation}
which obeys \cref{thm:response-Herglotz,thm:positive-response-kernel}.
\end{corollary}

\begin{corollary}[Response of the Hilbert block branch]
\label{cor:Hilbert-block-response}
Under \cref{ass:block-generator}, let \(\varepsilon=+1\), let
\(\Gamma_+\in\mathcal B(\scrE,\scrHp)\), and set
\begin{equation}
  \widetilde\Gamma_+h:=(\Gamma_+h,0)\in\scrK.
\end{equation}
Because \(\mathbb L^{(+1)}\) is Hilbert-self-adjoint, its response is
Herglotz and, whenever \(z\in\rho(L_-)\cap\rho(\mathbb L^{(+1)})\),
\begin{equation}
  \boxed{
  \widetilde\Gamma_+^*(\mathbb L^{(+1)}-z)^{-1}\widetilde\Gamma_+
  =\Gamma_+^*S_{+1}(z)^{-1}\Gamma_+.}
  \label{eq:Hilbert-block-response}
\end{equation}
\end{corollary}

\begin{remark}[Positivity boundary]
\label{rem:positivity-boundary}
The analogous Krein branch \(\mathbb L^{(-1)}\) is generally not
Hilbert-self-adjoint.  Its compressed resolvent therefore does not inherit
\cref{thm:response-Herglotz} without an additional positive realization or a
separate metric theorem.  No such positivity is assumed here.
\end{remark}

\subsection{Cayley transform and the Schur kernel}

The upper-half-plane Cayley transform converts a passive response into a
contractive transfer function.

\begin{definition}[Schur response]
\label{def:Schur-response}
For \(z\in\Cplus\), define
\begin{equation}
  \boxed{
  \Sresp(z)
  :=\bigl(\Mresp(z)-i\idop_{\scrE}\bigr)
    \bigl(\Mresp(z)+i\idop_{\scrE}\bigr)^{-1}.}
  \label{eq:Schur-response}
\end{equation}
\end{definition}

\begin{lemma}[Cayley denominator]
\label{lem:Cayley-denominator}
For every \(z\in\Cplus\), the operator
\(\Mresp(z)+i\idop_{\scrE}\) is boundedly invertible.
\end{lemma}

\begin{proof}
For \(h\in\scrE\),
\begin{equation}
  \|\bigl(\Mresp(z)+i\idop\bigr)h\|\,\|h\|
  \geq
  \operatorname{Im}\inner{h}{(\Mresp(z)+i\idop)h}
  \geq\|h\|^2.
\end{equation}
Thus the operator is bounded below.  For its adjoint,
\begin{equation}
  -\operatorname{Im}\inner{h}{(\Mresp(z)^*-i\idop)h}
  \ge \|h\|^2,
\end{equation}
so the adjoint is also bounded below and has trivial kernel.  Therefore the
original range is dense.  Since a bounded-below operator has closed range,
its range is all of \(\scrE\).
\end{proof}

\begin{theorem}[Schur contractivity]
\label{thm:Schur-contractivity}
The function \(\Sresp\) is analytic on \(\Cplus\) and satisfies
\begin{equation}
  \boxed{
  \idop_{\scrE}-\Sresp(z)^*\Sresp(z)
  =4\bigl(\Mresp(z)^*-i\idop\bigr)^{-1}
  \operatorname{Im}\Mresp(z)
  \bigl(\Mresp(z)+i\idop\bigr)^{-1}
  \geq0.}
  \label{eq:Schur-contractive-identity}
\end{equation}
In particular, \(\|\Sresp(z)\|\leq1\).
\end{theorem}

\begin{proof}
Put \(M=\Mresp(z)\).  Multiplying
\(\idop-\Sresp(z)^*\Sresp(z)\) on the left by \(M^*-i\idop\) and on the
right by \(M+i\idop\) gives
\begin{align}
  &(M^*-i\idop)(M+i\idop)
   -(M^*+i\idop)(M-i\idop)\\
  &\hspace{3cm}=2i(M^*-M)=4\operatorname{Im}M.
\end{align}
The asserted factorization and positivity follow.
\end{proof}

\begin{definition}[de Branges--Rovnyak kernel]
\label{def:dBR-kernel}
For \(z,w\in\Cplus\), define
\begin{equation}
  \KS(z,w)
  :=\frac{\idop_{\scrE}-\Sresp(z)\Sresp(w)^*}
  {-i(z-\overline w)}.
  \label{eq:dBR-kernel}
\end{equation}
\end{definition}

\begin{theorem}[Positive Schur defect kernel]
\label{thm:positive-Schur-kernel}
The kernel \(\KS\) is positive.  It is related to the Nevanlinna kernel by
\begin{equation}
  \boxed{
  \KS(z,w)
  =2\bigl(\Mresp(z)+i\idop\bigr)^{-1}
  \KM(z,w)
  \bigl(\Mresp(w)^*-i\idop\bigr)^{-1}.}
  \label{eq:Schur-kernel-factor}
\end{equation}
Consequently, \(\KS\) is the reproducing kernel of an operator-valued
de Branges--Rovnyak space.
\end{theorem}

\begin{proof}
A direct calculation gives
\begin{align}
 &(\Mresp(z)+i\idop)
  \bigl(\idop-\Sresp(z)\Sresp(w)^*\bigr)
  (\Mresp(w)^*-i\idop)\\
 &\hspace{3cm}=-2i\bigl(\Mresp(z)-\Mresp(w)^*\bigr).
\end{align}
Division by \(-i(z-\overline w)\) yields
\eqref{eq:Schur-kernel-factor}.  Positivity follows from
\cref{thm:positive-response-kernel} by congruence with the analytic factor
\(\sqrt2(\Mresp(z)+i\idop)^{-1}\).
\end{proof}

\subsection{Spectral measure and time-domain response}

Let \(E_H\) denote the projection-valued spectral measure of \(H\).

\begin{definition}[Compressed spectral measure]
\label{def:compressed-spectral-measure}
The positive operator-valued measure associated with \((H,\Gamma)\) is
\begin{equation}
  \Sigma_\Gamma(\Delta)
  :=\Gamma^*E_H(\Delta)\Gamma,
  \qquad \Delta\subseteq\RR\text{ Borel}.
  \label{eq:compressed-spectral-measure}
\end{equation}
Its total mass is
\begin{equation}
  \Sigma_\Gamma(\RR)=\Gamma^*\Gamma.
  \label{eq:spectral-measure-total-mass}
\end{equation}
\end{definition}

\begin{theorem}[Spectral and temporal representations]
\label{thm:spectral-temporal-response}
For \(z\in\CC\setminus\RR\),
\begin{equation}
  \boxed{
  \Mresp(z)=\int_{\RR}\frac{1}{\lambda-z}\,
  d\Sigma_\Gamma(\lambda).}
  \label{eq:response-spectral-representation}
\end{equation}
The time-domain compressed evolution
\begin{equation}
  C_\Gamma(t):=\Gamma^*e^{-itH}\Gamma
  \label{eq:time-response}
\end{equation}
satisfies
\begin{equation}
  \boxed{
  C_\Gamma(t)=\int_{\RR}e^{-it\lambda}\,
  d\Sigma_\Gamma(\lambda).}
  \label{eq:time-response-spectral}
\end{equation}
Moreover, \(C_\Gamma\) is operator-valued positive definite:
\begin{equation}
  \sum_{i,j=1}^{N}
  \inner{h_i}{C_\Gamma(t_i-t_j)h_j}_{\scrE}\geq0.
  \label{eq:time-response-positive-definite}
\end{equation}
\end{theorem}

\begin{proof}
The first two formulas are direct applications of the spectral theorem.
For positive definiteness, the left-hand side of
\eqref{eq:time-response-positive-definite} is
\begin{equation}
  \left\|
  \sum_{j=1}^{N}e^{it_jH}\Gamma h_j
  \right\|_{\scrX}^{2}.
\end{equation}
\end{proof}

\subsection{Minimal response realizations}

Only the spectral subspace reached from the port channel can influence the
response.

\begin{definition}[Response-generated subspace]
\label{def:response-generated-subspace}
Define
\begin{equation}
  \Hmin
  :=\overline{\operatorname{span}}
  \left\{
    R_H(z)\Gamma h:
    z\in\CC\setminus\RR,\ h\in\scrE
  \right\}
  \subseteq\scrX.
  \label{eq:minimal-subspace}
\end{equation}
The realization \((H,\Gamma)\) is called minimal if
\(\Hmin=\scrX\).
\end{definition}

\begin{theorem}[Canonical minimal reduction]
\label{thm:canonical-minimal-reduction}
The subspace \(\Hmin\) reduces \(H\), contains \(\Ran\Gamma\), and the
restricted operator
\begin{equation}
  H_{\min}:=\restr{H}{\Dom H\cap\Hmin}
  \label{eq:Hmin}
\end{equation}
is self-adjoint on \(\Hmin\).  Viewing \(\Gamma\) as a map into \(\Hmin\),
\begin{equation}
  \boxed{
  \Gamma^*(H-z)^{-1}\Gamma
  =\Gamma^*(H_{\min}-z)^{-1}\Gamma.}
  \label{eq:minimal-response-invariance}
\end{equation}
Thus every realization admits a canonical minimal part with exactly the same
response.
\end{theorem}

\begin{proof}
For \(z,\zeta\in\CC\setminus\RR\), the resolvent identity gives
\begin{equation}
  R_H(\zeta)R_H(z)
  =\frac{R_H(\zeta)-R_H(z)}{\zeta-z}.
\end{equation}
Hence \(\Hmin\) is invariant under every resolvent.  It is also invariant
under their adjoints because \(R_H(\zeta)^*=R_H(\overline\zeta)\).  It
therefore reduces \(H\).  Furthermore,
\begin{equation}
  -iy(H-iy)^{-1}\longrightarrow\idop
  \quad\text{strongly as }y\to+\infty,
\end{equation}
so \(\Ran\Gamma\subseteq\Hmin\).  The remaining statements follow from the
spectral theorem for reducing subspaces.
\end{proof}

\begin{theorem}[Unitary uniqueness of minimal realizations]
\label{thm:minimal-uniqueness}
Let \((H_1,\Gamma_1)\) on \(\scrX_1\) and \((H_2,\Gamma_2)\) on
\(\scrX_2\) be minimal realizations over the same port space \(\scrE\).  If
\begin{equation}
  \Gamma_1^*(H_1-z)^{-1}\Gamma_1
  =\Gamma_2^*(H_2-z)^{-1}\Gamma_2
  \label{eq:equal-minimal-responses}
\end{equation}
for all \(z\in\CC\setminus\RR\), then there is a unique unitary
\(U:\scrX_1\to\scrX_2\) satisfying
\begin{equation}
  UH_1=H_2U,
  \qquad
  U\Gamma_1=\Gamma_2.
  \label{eq:minimal-unitary-equivalence}
\end{equation}
\end{theorem}

\begin{proof}
Let \(R_j(z)=(H_j-z)^{-1}\).  On finite response sums define
\begin{equation}
  U_0\left(\sum_{k=1}^N R_1(z_k)\Gamma_1h_k\right)
  :=\sum_{k=1}^N R_2(z_k)\Gamma_2h_k.
  \label{eq:U0-response-sums}
\end{equation}
For \(z,w\in\CC\setminus\RR\) with \(w\ne\overline z\), the
resolvent identity gives
\begin{equation}
  \inner{R_j(z)\Gamma_jh}{R_j(w)\Gamma_jk}_{\scrX_j}
  =\inner{h}{
    \frac{M_j(w)-M_j(z)^*}{w-\overline z}
    k}_{\scrE}.
  \label{eq:minimal-Gram-formula}
\end{equation}
The case \(w=\overline z\) follows by norm-continuous limiting from nearby
nonreal parameters.  The common response therefore gives identical Gram matrices for arbitrary
finite response sums in the two realizations.  Equation \eqref{eq:U0-response-sums} is consequently
well-defined and isometric.  Minimality makes its initial and final spaces
dense, so it extends uniquely to a unitary \(U\).

The resolvent identity shows on the dense response span that
\begin{equation}
  UR_1(\zeta)=R_2(\zeta)U,
  \qquad \zeta\in\CC\setminus\RR.
\end{equation}
Hence \(UH_1=H_2U\).  Finally,
\begin{equation}
  -iyR_j(iy)\Gamma_j\longrightarrow\Gamma_j
  \qquad (y\to+\infty)
\end{equation}
strongly, which yields \(U\Gamma_1=\Gamma_2\).
\end{proof}

\begin{corollary}[Finite response dimension and indirect directions]
\label{cor:finite-response-defect}
Suppose the compressed spectral measure has finite support at distinct
real numbers \(\lambda_1,\ldots,\lambda_m\), with finite-rank weights
\(R_j\succeq0\), so that
\begin{equation}
  \Mresp(z)=\sum_{j=1}^{m}\frac{R_j}{\lambda_j-z}.
  \label{eq:finite-response-weights}
\end{equation}
The space reached by the dynamics outside the direct port range is
\begin{equation}
  \mathcal D_\Gamma
  :=\Hmin\ominus\overline{\Ran\Gamma}
  =\Hmin\cap\Ker\Gamma^*.
  \label{eq:indirect-response-space}
\end{equation}
Its dimension is determined by the response:
\begin{equation}
  \boxed{
  \dim\mathcal D_\Gamma
  =\sum_{j=1}^{m}\operatorname{rank}R_j
   -\operatorname{rank}\!\left(\sum_{j=1}^{m}R_j\right).}
  \label{eq:finite-response-defect-count}
\end{equation}
\end{corollary}

\begin{proof}
Let \(E_j=E_H(\{\lambda_j\})\).  The spectral measure gives
\(R_j=\Gamma^*E_j\Gamma=(E_j\Gamma)^*(E_j\Gamma)\), hence
\(\dim\Ran(E_j\Gamma)=\operatorname{rank}R_j\).
The response-generated space is the orthogonal direct sum of these ranges.
Indeed, the finite spectral support makes their sum invariant under every
resolvent, and the distinct spectral projections can be recovered on that
sum by polynomial interpolation in \(H\).  Thus
\(\dim\Hmin=\sum_j\operatorname{rank}R_j\).
Also \(\Gamma^*\Gamma=\sum_jR_j\), so
\(\dim\Ran\Gamma=\operatorname{rank}(\sum_jR_j)\).
Both spaces are finite-dimensional and the direct range lies in \(\Hmin\).
Subtracting their dimensions proves the formula.
\end{proof}

After restriction to the finite-dimensional active port, the identity
$\dim\Hmin=\sum_j\operatorname{rank}R_j$ is Gilbert's
simple-pole minimal-realization formula applied to a positive spectral
measure \cite[Theorem~7]{Gilbert1963}.  The final-defect count
subtracts the dimension of the direct preparation range from that
minimal dimension.

The count concerns a specified port and its complete response.  An
invertible change of input coordinates preserves it.  Positivity gives
\(\Ran(\sum_j R_j)=\sum_j\Ran R_j\), since the two finite-dimensional
spaces have orthogonal complement \(\bigcap_j\Ker R_j\).
Thus the count is the dimension of the space of relations
\(\sum_j x_j=0\) with \(x_j\in\Ran R_j\).  It vanishes precisely when
those ranges form a direct sum in the port space.  Adding independent
ports can change the count.

If a realization takes \(B=\Gamma\) with codomain \(\Hmin\), then
\(\Ker B^*=\mathcal D_\Gamma\).  With the full codomain \(\scrX\), one
has instead
\begin{equation}
  \Ker\Gamma^*=\mathcal D_\Gamma\oplus(\scrX\ominus\Hmin).
  \label{eq:ambient-response-defect}
\end{equation}
The second summand is orthogonal to every state generated by the port
and is undetermined by the response.  Injectivity of the input channel
requires \(\Ker\Gamma=\Ker(\sum_j R_j)=\{0\}\).
The choice of transition channel and any ground correspondence are
additional realization data, as described in \cref{eq:realization-datum}.

\begin{proposition}[Permanent invisibility to a fixed port]
\label{prop:permanent-port-invisibility}
For the self-adjoint realization $(H,\Gamma)$, define
\begin{equation}
 \mathcal N_\Gamma
 =\{x\in\scrX:\Gamma^*e^{-itH}x=0
                 \text{ for every }t\in\RR\}.
\end{equation}
Then $\mathcal N_\Gamma=\Hmin^\perp$.  In particular, no nonzero
vector in $\mathcal D_\Gamma$ is permanently invisible to the same
port under the full unitary group.
\end{proposition}
\begin{proof}
By the spectral theorem, the closed span of
$e^{itH}\Gamma h$, with $t\in\RR$ and $h\in\scrE$, is the same
reducing subspace as the resolvent span in
\eqref{eq:minimal-subspace}.  For instance, Laplace transforms of
the unitary group give its nonreal resolvents, and invariance under
all resolvents gives the spectral projections and hence the group.
The displayed condition says exactly that $x$ is orthogonal to
this span.
\end{proof}

In finite dimensions, this is the controllability--observability
duality for a system with input map $\Gamma$ and output map
$\Gamma^*$ \cite[Sec.~5]{Kalman1963}.  The spectral-theorem argument
above establishes the stated version for a possibly
infinite-dimensional self-adjoint realization.

The conclusion concerns the stated port and all real times.  It
does not assert detectability in every finite observation window
or for every other observable.  A dynamically accessible final
defect and a permanently decoupled ambient sector are therefore
different objects.


\subsection{Symmetry-resolved response and departure from a port}
\label{sec:symmetry-response}

A conserved generator can preserve a reference state while separating the
energies present in a prepared superposition.  Commutation of the generator
with the Hamiltonian does not imply that the direct port range is invariant.
The following statement uses the same time convention as
\eqref{eq:time-response}, namely $\hbar=1$.

\begin{proposition}[Opposite symmetry sectors]
\label{prop:opposite-sector-response}
Let $H_0$ and $L$ be strongly commuting self-adjoint operators on $\scrX$,
and define $H_\nu=H_0-\nu L$ by their joint functional calculus.  Suppose
that a normalized vector $g$ satisfies $H_0g=E_0g$ and $Lg=0$.  Let
$\ell>0$, and let $q_\pm\in\Ker(L\mp\ell\idop)$ be normalized vectors
with the same spectral measure for $H_0-E_0$:
\begin{equation}
 \inner{q_+}{E_{H_0-E_0}(\Delta)q_+}
 =\inner{q_-}{E_{H_0-E_0}(\Delta)q_-}
 =:\mu(\Delta).
 \label{eq:opposite-sector-measure}
\end{equation}
Set
\begin{equation}
 e:=\frac{q_++q_-}{\sqrt2},\qquad
 d:=\frac{q_+-q_-}{i\sqrt2},\qquad
 \Gamma(c_0,c_1):=c_0g+c_1e,
 \qquad C_0(t):=\int_{\RR}e^{-it\lambda}\,d\mu(\lambda).
 \label{eq:opposite-sector-port}
\end{equation}
Then $g,e,d$ are orthonormal, $\Gamma$ is an isometry, and the normalized
excitation response is
\begin{equation}
 C_\nu(t):=\inner{e}{e^{-it(H_\nu-E_0)}e}
 =C_0(t)\cos(\nu\ell t).
 \label{eq:opposite-sector-factorization}
\end{equation}
Writing $\psi(t)=e^{-it(H_\nu-E_0)}e$, the weight outside the direct port
range equals
\begin{equation}
 \|P_{\mathcal D_\Gamma}\psi(t)\|^2
 =1-|C_0(t)|^2\cos^2(\nu\ell t),
 \label{eq:opposite-sector-departure}
\end{equation}
where $\mathcal D_\Gamma$ is defined using the minimal response space of
$(H_\nu,\Gamma)$.  For $\nu\ne0$, this weight equals one at
$t_* =\pi/(2|\nu|\ell)$.  The weight on the fixed analysis line $\CC d$ is
\begin{equation}
 |\inner{d}{\psi(t)}|^2
 =|C_0(t)|^2\sin^2(\nu\ell t),
 \label{eq:opposite-sector-fixed-line}
\end{equation}
and the weight orthogonal to $\Span\{g,e,d\}$ is $1-|C_0(t)|^2$.
\end{proposition}

\begin{proof}
Distinct eigenspaces of $L$ are orthogonal.  Strong commutation implies
\[
 e^{-it(H_\nu-E_0)}q_\pm
 =e^{\pm it\nu\ell}e^{-it(H_0-E_0)}q_\pm.
\]
Taking the inner products with $e$ and $d$ gives respectively
$C_0(t)\cos(\nu\ell t)$ and $-C_0(t)\sin(\nu\ell t)$.
The evolving vector remains orthogonal to $g$ and belongs to $\Hmin$.
Consequently its projection onto $\Ran\Gamma$ has squared norm
$|C_\nu(t)|^2$, which proves \eqref{eq:opposite-sector-departure}.
The last identity follows by subtracting the squared projections onto $e$
and $d$ from the conserved unit norm.
\end{proof}

Matched measures in \eqref{eq:opposite-sector-measure} can follow from a
unitary reflection, or an antiunitary symmetry such as time reversal,
that commutes with $H_0$, reverses $L$, and exchanges $q_+$ and $q_-$.
In the antiunitary case the spectral projections are preserved and their
expectations are real, so the two spectral measures still agree.
The fixed-reference identity $Lg=0$ by itself proves
that $g$ remains an eigenvector, not that it remains the lowest-energy
state: preservation of a ground requires a separate no-crossing condition.
Similarly, complete departure from the direct port does not ensure a
complete return; at $\nu\ell t\in\pi\ZZ$, the return weight is
$|C_0(t)|^2$.

The line $\CC d$ is an analysis direction in the ambient carrier.  If the
two shifted spectral branches overlap, $d$ need not belong to $\Hmin$.
It must not then be identified with an individually reachable minimal
defect line.  Formula \eqref{eq:opposite-sector-departure} still concerns
the response-generated defect because $\psi(t)\in\Hmin$.  An isolated
pair of distinct eigenvalues gives the simpler finite case below.


\subsection{Heat-smoothed singular inputs}

A point trace or boundary delta is usually not a vector of the carrier
Hilbert space.  Gaussian functional calculus gives a canonical way to turn
such formal data into a bounded port channel.

For \(s\geq0\), let
\begin{equation}
  \scrX_s:=\Dom(\idop+H^2)^{s/2}
\end{equation}
with its graph norm, and let \(\scrX_{-s}\) denote its anti-dual with pivot
\(\scrX\).

\begin{assumption}[Distributional trace]
\label{ass:distributional-trace}
Let
\begin{equation}
  \gamma:\scrX_s\longrightarrow\scrE
\end{equation}
be bounded.  Its adjoint is interpreted as
\begin{equation}
  \gamma^*:\scrE\longrightarrow\scrX_{-s}.
\end{equation}
\end{assumption}

\begin{theorem}[Heat-smoothed distributional port]
\label{thm:heat-smoothed-distributional-port}
For every \(\tau>0\), the operator
\begin{equation}
  \boxed{
  \Gamma_{\tau,\gamma}
  :=e^{-\tau H^2/2}\gamma^*
  \in\mathcal B(\scrE,\scrX)}
  \label{eq:heat-smoothed-distributional-port}
\end{equation}
is bounded.  Its compressed resolvent
\begin{equation}
  M_{\tau,\gamma}(z)
  :=\Gamma_{\tau,\gamma}^*(H-z)^{-1}
  \Gamma_{\tau,\gamma}
  \label{eq:heat-smoothed-distributional-response}
\end{equation}
is therefore a Herglotz response with finite total spectral mass
\begin{equation}
  \Sigma_{\tau,\gamma}(\RR)
  =\Gamma_{\tau,\gamma}^*\Gamma_{\tau,\gamma}.
\end{equation}
\end{theorem}

\begin{proof}
The spectral multiplier
\begin{equation}
  (1+\lambda^2)^{s/2}e^{-\tau\lambda^2/2}
\end{equation}
is bounded on \(\RR\).  Hence \(e^{-\tau H^2/2}\) extends boundedly from
\(\scrX_{-s}\) to \(\scrX\).  Composition with \(\gamma^*\) proves
\eqref{eq:heat-smoothed-distributional-port}; the response claims follow from
\cref{thm:response-Herglotz,def:compressed-spectral-measure}.
\end{proof}

\begin{assumption}[Discrete spectral realization]
\label{ass:discrete-response-spectrum}
Assume that \(H\) has a complete orthonormal eigenbasis
\(\{e_n\}_{n\in I}\) with
\begin{equation}
  He_n=\lambda_ne_n,
  \qquad \lambda_n\in\RR.
  \label{eq:discrete-spectrum}
\end{equation}
Let \(\{\gamma_n\}_{n\in I}\subseteq\scrE\) specify a formal port through
its spectral coefficients.
\end{assumption}

\begin{theorem}[Heat-smoothed port channel]
\label{thm:heat-smoothed-channel}
Fix \(\tau>0\) and suppose
\begin{equation}
  \sum_{n\in I}e^{-\tau\lambda_n^2}\|\gamma_n\|_{\scrE}^2<\infty.
  \label{eq:heat-smoothing-summability}
\end{equation}
Then
\begin{equation}
  \Gamma_\tau h
  :=\sum_{n\in I}
  e^{-\tau\lambda_n^2/2}
  \inner{\gamma_n}{h}_{\scrE}\,e_n
  \label{eq:heat-smoothed-input}
\end{equation}
defines a Hilbert--Schmidt operator \(\Gamma_\tau:\scrE\to\scrX\), and
its response is
\begin{equation}
  \boxed{
  M_\tau(z)
  =\sum_{n\in I}
  \frac{e^{-\tau\lambda_n^2}}
  {\lambda_n-z}\,
  \gamma_n\gamma_n^*.}
  \label{eq:heat-smoothed-response}
\end{equation}
The series converges locally uniformly in operator norm on
\(\CC\setminus\RR\).  The total compressed spectral mass is
\begin{equation}
  \boxed{
  \Gamma_\tau^*\Gamma_\tau
  =\sum_{n\in I}e^{-\tau\lambda_n^2}\gamma_n\gamma_n^*,}
  \label{eq:heat-smoothed-total-mass}
\end{equation}
with convergence in trace norm.
\end{theorem}

\begin{proof}
The Hilbert--Schmidt norm is exactly the sum in
\eqref{eq:heat-smoothing-summability}.  On every compact set \(K\subset\CC\setminus\RR\),
\(\sup_{z\in K}|\lambda_n-z|^{-1}\leq\operatorname{dist}(K,\RR)^{-1}\).
Hence the response series converges locally uniformly in operator norm.
The total-mass identity follows by direct computation, and the trace norm of
each positive rank-one summand is
\(e^{-\tau\lambda_n^2}\|\gamma_n\|^2\).
\end{proof}

\begin{remark}[Normalization]
\label{rem:heat-channel-normalization}
For a scalar port space, if
\begin{equation}
  c_\tau:=\|\Gamma_\tau 1\|^2>0,
\end{equation}
then \(\widehat\Gamma_\tau:=c_\tau^{-1/2}\Gamma_\tau\) has a compressed
spectral measure of total mass one.  For a general port space, write the polar decomposition
\(\Gamma_\tau=V_\tau(\Gamma_\tau^*\Gamma_\tau)^{1/2}\).  The partial
isometry \(V_\tau\), restricted to the support of
\(\Gamma_\tau^*\Gamma_\tau\), is the bounded support-normalized port.  This
formulation remains valid even when the positive operator has spectrum
accumulating at zero.
\end{remark}
% ===== END INLINED FILE: sections/08_boundary_response.tex =====


% ===== BEGIN INLINED FILE: sections/09_examples.tex =====
\section{Canonical realizations}
\label{sec:examples}

The transition algebra fixes the relations among
\(\Aop_B,\Vop_B,\DB_B,\KB_B\), and \(\DeltaB_B\), but it does not fix the
spectrum.  Spectral type enters through the closed channel \(B\).  The
following examples exhibit finite, continuous, oscillator, and compact
quasiperiodic realizations of the same algebra.

\subsection{A scalar channel}

\begin{example}[One-dimensional realization]
\label{ex:scalar-channel}
Let \(\scrHp=\scrHm=\CC\) and let \(Bx=bx\) for \(b\in\CC\).  Then
\begin{align}
  \DB_B
  &=\begin{pmatrix}0&\overline b\\ b&0\end{pmatrix},
  &
  \KB_B
  &=\begin{pmatrix}0&\overline b\\ -b&0\end{pmatrix},
  \label{eq:scalar-odd-operators}\\
  \DeltaB_B
  &=|b|^2\idop_{\CC^2}.
  \label{eq:scalar-laplacian}
\end{align}
If \(b\ne0\), the active support is all of \(\CC^2\), the polar phase is
\(u_B=b/|b|\), and \(\ind B=0\).  Taking
\begin{equation}
  H_b:=|b|^2\quad\text{on }\CC,
  \qquad
  \Gamma_b:=\idop_\CC,
\end{equation}
gives the minimal scalar response
\begin{equation}
  M_b(z)=\frac{1}{|b|^2-z}.
  \label{eq:scalar-response}
\end{equation}
\end{example}

This example contains no spectral multiplicity or hidden domain issue.  It is
the local normal form of a single active transition amplitude.

\subsection{An invariant ground with a one-dimensional response defect}
\label{sec:three-state-response}

\begin{example}[Split excitations and a grounded two-dimensional port]
\label{ex:split-excitation-response}
Let $g,f_+,f_-$ be an orthonormal basis of $\scrX=\CC^3$ and take
\begin{equation}
 H_0g=E_0g,\qquad H_0f_\pm=E_2f_\pm,
 \qquad Lg=0,\qquad Lf_\pm=\pm\ell f_\pm,
 \label{eq:three-state-data}
\end{equation}
where $\ell>0$ and $0<|\nu|\ell<E_2-E_0$.  For
$H_\nu=H_0-\nu L$, $g$ is therefore the unique ground and the two
excitation energies are $E_2\mp\nu\ell$.  Define
\begin{equation}
 e=\frac{f_++f_-}{\sqrt2},\qquad
 d=\frac{f_+-f_-}{i\sqrt2},\qquad
 \Gamma(c_0,c_1)=c_0g+c_1e.
 \label{eq:three-state-port}
\end{equation}
In the standard basis of $\scrE=\CC^2$, the complete response is
\begin{equation}
 M_\Gamma(z)=
 \begin{pmatrix}
 (E_0-z)^{-1}&0\\
 0&\dfrac{1}{2(E_2-\nu\ell-z)}
    +\dfrac{1}{2(E_2+\nu\ell-z)}
 \end{pmatrix}.
 \label{eq:three-state-response}
\end{equation}
Its three distinct atoms have weights
\begin{equation}
 R_0=\begin{pmatrix}1&0\\0&0\end{pmatrix},\qquad
 R_+=R_-=\frac12\begin{pmatrix}0&0\\0&1\end{pmatrix},
 \qquad R_0+R_++R_-=\idop_{\CC^2}.
 \label{eq:three-state-weights}
\end{equation}
Thus \cref{cor:finite-response-defect} gives
\begin{equation}
 \Hmin=\scrX,\qquad \dim\mathcal D_\Gamma=3-2=1,
 \qquad \mathcal D_\Gamma=\CC d.
 \label{eq:three-state-minimal-defect}
\end{equation}
Explicitly,
\begin{equation}
 e^{-it(H_\nu-E_0)}e
 =e^{-it(E_2-E_0)}
   \bigl[\cos(\nu\ell t)e-\sin(\nu\ell t)d\bigr].
 \label{eq:three-state-transfer}
\end{equation}
The defect weight is $\sin^2(\nu\ell t)$, and both positive response
weights are $1/2$.  At zero bias, the two excited atoms coincide and must
be merged: the minimal carrier drops to $\Span\{g,e\}$ and the
response-generated defect vanishes.
\end{example}

Taking $\scrHp=\CC^2$, $\scrHm=\Hmin$ and $B=\Gamma$ yields
$B^*B=\idop_{\CC^2}$, $\Ker B=\{0\}$ and $\Ker B^*=\CC d$.
The normalized port ground $(1,0)$ maps to $g$.  If its input Hamiltonian
is the compression $\Gamma^*H_\nu\Gamma=\diag(E_0,E_2)$, the input
and output grounds correspond exactly.  The response Hamiltonian
$H_\nu$ is additional realization data: it is neither $B^*B$ nor $BB^*$.
The port is not dynamically invariant when $\nu\ne0$, despite
$[H_\nu,L]=0$.


\subsection{The translation channel and continuous spectrum}

\begin{example}[Free translation channel]
\label{ex:translation-channel}
Let
\begin{equation}
  \scrHp=\scrHm=L^2(\RR),
  \qquad
  B=D_x:=-i\frac{d}{dx},
  \qquad
  \Dom D_x=H^1(\RR).
  \label{eq:translation-channel}
\end{equation}
Then \(D_x=D_x^*\), and the partner operators coincide:
\begin{equation}
  B^*B=BB^*=D_x^2=-\frac{d^2}{dx^2},
  \qquad
  \Dom D_x^2=H^2(\RR).
  \label{eq:free-partner}
\end{equation}
Under the unitary Fourier transform, \(D_x\) is multiplication by \(k\) and
\(D_x^2\) is multiplication by \(k^2\).  Hence
\begin{equation}
  \spec(D_x^2)=[0,\infty)
  \label{eq:free-spectrum}
\end{equation}
is purely continuous.
\end{example}

The channel has trivial kernel but no support gap.  Indeed, choose normalized
vectors \(f_n\) whose Fourier transforms are supported in
\([-1/n,1/n]\).  Then
\begin{equation}
  \|D_xf_n\|\le\frac1n,
  \qquad
  \|f_n\|=1,
  \label{eq:no-translation-gap}
\end{equation}
so no positive lower bound can hold on \((\Ker D_x)^\perp=L^2(\RR)\).  By the
closed-range criterion of \cref{thm:closed-range-criterion}, \(\Ran D_x\) is not
closed.

For any normalized \(g\in L^2(\RR)\), let
\(\Gamma_gc=cg\) and take \(H_{\mathrm{free}}=D_x^2\).  The corresponding
response is
\begin{equation}
  M_g(z)
  =\int_{\RR}\frac{|\widehat g(k)|^2}{k^2-z}\,dk,
  \qquad z\in\CC\setminus[0,\infty),
  \label{eq:free-response}
\end{equation}
where \(\widehat g\) is the unitary Fourier transform of \(g\).  Thus an
absolutely continuous response measure and the absence of a support gap are
compatible with the same A/V algebra used in the discrete examples below.

\subsection{The oscillator channel and the integer heat trace}

Let \(\mathfrak h=\ell^2(\NN_0)\) with orthonormal basis
\((e_n)_{n\ge0}\).  Define
\begin{align}
  ae_0&=0,
  &ae_n&=\sqrt n\,e_{n-1}\quad(n\ge1),
  \label{eq:annihilation-action}\\
  a^*e_n&=\sqrt{n+1}\,e_{n+1}.
  \label{eq:creation-action}
\end{align}
Their natural domains are
\begin{align}
  \Dom a
  &=\left\{x=(x_n):\sum_{n\ge0}n|x_n|^2<\infty\right\},
  \label{eq:annihilation-domain}\\
  \Dom a^*
  &=\left\{x=(x_n):\sum_{n\ge0}(n+1)|x_n|^2<\infty\right\}.
  \label{eq:creation-domain}
\end{align}
Let \(N=a^*a\) be the number operator.

\begin{theorem}[Oscillator transition channel]
\label{thm:oscillator-channel}
Take \(\scrHp=\scrHm=\mathfrak h\) and
\begin{equation}
  B=a^*.
  \label{eq:oscillator-B}
\end{equation}
Then
\begin{equation}
  B^*B=aa^*=N+\idop,
  \qquad
  BB^*=a^*a=N.
  \label{eq:oscillator-partners}
\end{equation}
The range of \(B\) is closed,
\begin{equation}
  \Ker B=\{0\},
  \qquad
  \Ker B^*=\operatorname{span}\{e_0\},
  \qquad
  \ind B=-1.
  \label{eq:oscillator-index}
\end{equation}
The positive return operator has spectrum
\begin{equation}
  \spec(B^*B)=\{1,2,3,\ldots\},
  \label{eq:positive-integer-spectrum}
\end{equation}
whereas its partner contains one additional zero mode.
\end{theorem}

\begin{proof}
The identities in \eqref{eq:oscillator-partners} follow on the basis vectors.
Moreover,
\begin{equation}
  \|a^*x\|^2
  =\sum_{n\ge0}(n+1)|x_n|^2
  \ge\|x\|^2,
\end{equation}
so \(a^*\) is bounded below and has closed range.  The kernel statements are
immediate, and the index follows from \cref{thm:index-formulas}.  The
spectral statement is the diagonal action of \(N+\idop\).
\end{proof}

\begin{corollary}[Index defect in the heat trace]
\label{cor:oscillator-heat-index}
For every \(t>0\),
\begin{equation}
  \Tr\bigl(e^{-t(N+\idop)}\bigr)
  -\Tr\bigl(e^{-tN}\bigr)
  =-1
  =\ind B.
  \label{eq:oscillator-linear-heat-index}
\end{equation}
The same unpaired zero mode gives
\begin{equation}
  \Tr\bigl(e^{-\pi t(N+\idop)^2}\bigr)
  -\Tr\bigl(e^{-\pi tN^2}\bigr)
  =-1.
  \label{eq:oscillator-square-heat-index}
\end{equation}
\end{corollary}

\begin{theorem}[Mellin transform of the positive integer heat trace]
\label{thm:oscillator-zeta}
For \(\operatorname{Re}s>1\),
\begin{align}
  \int_0^\infty
  t^{s/2-1}
  \Tr\bigl(e^{-\pi t(B^*B)^2}\bigr)\,dt
  &=\int_0^\infty
  t^{s/2-1}\sum_{n=1}^\infty e^{-\pi n^2t}\,dt
  \label{eq:oscillator-Mellin-start}\\
  &=\pi^{-s/2}\Gamma\!\left(\frac s2\right)\zeta(s).
  \label{eq:oscillator-Mellin-zeta}
\end{align}
\end{theorem}

\begin{proof}
For \(\operatorname{Re}s>1\), positivity and absolute convergence permit
termwise integration.  The substitution \(u=\pi n^2t\) gives
\begin{equation}
  \int_0^\infty t^{s/2-1}e^{-\pi n^2t}\,dt
  =\pi^{-s/2}n^{-s}\Gamma\!\left(\frac s2\right).
\end{equation}
Summation over \(n\ge1\) proves the result.
\end{proof}

\begin{remark}
The right-hand side of \eqref{eq:oscillator-Mellin-zeta} is the standard
Gamma-completed zeta factor
\(\pi^{-s/2}\Gamma(s/2)\zeta(s)\).  It is not the entire Riemann
\(\xi\)-function, which includes the additional factor
\(\tfrac12s(s-1)\), and it makes no assertion about the spectrum of the
nontrivial zeta zeros.
\end{remark}

The oscillator also gives a simple minimal response.  For \(0<r<1\), set
\begin{equation}
  g_r:=\sqrt{1-r^2}\sum_{n=0}^\infty r^n e_n,
  \qquad
  \Gamma_rc:=cg_r.
  \label{eq:oscillator-cyclic-vector}
\end{equation}
Then \(\|g_r\|=1\), and for \(H_{\mathrm{osc}}=N+\idop\),
\begin{equation}
  M_r(z)
  =(1-r^2)\sum_{n=0}^\infty
  \frac{r^{2n}}{n+1-z}.
  \label{eq:oscillator-response}
\end{equation}
Every spectral coefficient of \(g_r\) is nonzero.  Since the spectral
projections of the simple pure-point operator \(N+\idop\) are recovered from
its resolvent, the resolvent orbit of \(g_r\) contains every basis direction;
hence the realization is minimal.  Thus the response uniquely determines
this oscillator realization
up to unitary equivalence by \cref{thm:minimal-uniqueness}.

\subsection{A quasiperiodic compact channel and a smoothed point input}

\begin{example}[First-order generator on a compact interval]
\label{ex:quasiperiodic-channel}
Let \(L>0\), \(\vartheta\in[0,2\pi)\), and
\begin{equation}
  \mathfrak h_L:=L^2(-L,L).
\end{equation}
Define
\begin{align}
  D_{L,\vartheta}
  &:=-i\frac{d}{du},
  \label{eq:quasiperiodic-generator}\\
  \Dom D_{L,\vartheta}
  &:=\left\{f\in H^1(-L,L):
  f(L)=e^{i\vartheta}f(-L)\right\}.
  \label{eq:quasiperiodic-domain}
\end{align}
Then \(D_{L,\vartheta}\) is self-adjoint with eigenvalues and normalized
eigenfunctions
\begin{align}
  \lambda_n
  &=\frac{2\pi n+\vartheta}{2L},
  \label{eq:quasiperiodic-eigenvalues}\\
  e_n(u)
  &=\frac{1}{\sqrt{2L}}
  \exp\bigl(i\lambda_n(u+L)\bigr),
  \qquad n\in\ZZ.
  \label{eq:quasiperiodic-eigenfunctions}
\end{align}
Taking \(B=D_{L,\vartheta}\) between two copies of \(\mathfrak h_L\) gives
\begin{equation}
  B^*B=BB^*=D_{L,\vartheta}^2.
  \label{eq:quasiperiodic-partners}
\end{equation}
The channel has index zero.  Its spectrum is discrete.  After removing a
possible zero mode, its active support has a positive gap from zero.
\end{example}

For \(u_0\in(-L,L)\), the point source \(\delta_{u_0}\) is not an element of
\(L^2(-L,L)\).  Regarded as an element of a negative Sobolev scale, its spectral coefficient
with respect to the point-evaluation trace is \(e_n(u_0)\), and hence
\begin{equation}
  |e_n(u_0)|^2=\frac{1}{2L}.
  \label{eq:uniform-point-weights}
\end{equation}
Consequently, the unsmoothed total spectral mass diverges:
\begin{equation}
  \sum_{n\in\ZZ}|e_n(u_0)|^2=\infty.
  \label{eq:unsmoothed-point-mass}
\end{equation}

For \(\tau>0\), apply \cref{thm:heat-smoothed-distributional-port} to the self-adjoint
generator \(H=D_{L,\vartheta}\).  The heat-smoothed point port vector is
\begin{equation}
  b_{\tau,u_0}
  :=e^{-\tau D_{L,\vartheta}^2/2}\delta_{u_0}
  =\sum_{n\in\ZZ}
  e^{-\tau\lambda_n^2/2}\overline{e_n(u_0)}e_n,
  \label{eq:smoothed-point-vector}
\end{equation}
with finite mass
\begin{equation}
  Z_{\tau,L,\vartheta}
  :=\|b_{\tau,u_0}\|^2
  =\frac{1}{2L}\sum_{n\in\ZZ}e^{-\tau\lambda_n^2}<\infty.
  \label{eq:theta-mass}
\end{equation}
The value is independent of \(u_0\).  The associated scalar response is
\begin{equation}
  M_{\tau,u_0}(z)
  =\frac{1}{2L}\sum_{n\in\ZZ}
  \frac{e^{-\tau\lambda_n^2}}{\lambda_n-z}.
  \label{eq:quasiperiodic-smoothed-response}
\end{equation}
After normalizing
\begin{equation}
  \widehat b_{\tau,u_0}
  :=Z_{\tau,L,\vartheta}^{-1/2}b_{\tau,u_0},
  \label{eq:normalized-smoothed-point}
\end{equation}
the compressed spectral measure has total mass one, while the internal
spectrum \((\lambda_n)\) remains unchanged.

\begin{remark}
When \(\vartheta=0\), the constant mode is a zero mode.  On its orthogonal
complement, the first nonzero magnitude is \(\pi/L\).  When
\(\vartheta\ne0\), the first-order gap is
\begin{equation}
  \frac{\operatorname{dist}(\vartheta,2\pi\ZZ)}{2L},
\end{equation}
and the return operator \(D_{L,\vartheta}^2\) has the square of this gap.
\end{remark}

\subsection{What is universal and what depends on the channel}

\begin{center}
\small
\begin{tabularx}{\textwidth}{>{\raggedright\arraybackslash}p{0.20\textwidth} >{\raggedright\arraybackslash}p{0.21\textwidth} >{\raggedright\arraybackslash}p{0.20\textwidth} X}
\toprule
\textbf{Channel} & \textbf{Positive return} & \textbf{Spectral type} & \textbf{Distinguishing feature} \\
\midrule
\(B=b\) on \(\CC\) & \(|b|^2\) & finite, one level & Full support for \(b\ne0\); index zero. \\
\(-i\,d/dx\) on \(\RR\) & \(-d^2/dx^2\) & continuous \([0,\infty)\) & No closed range and no support gap. \\
\(B=a^*\) & \(N+\idop\) & positive integers & Index \(-1\); exact theta/zeta heat trace. \\
\(D_{L,\vartheta}\) & \(D_{L,\vartheta}^2\) & discrete quasiperiodic & Point input becomes finite after Gaussian smoothing. \\
\bottomrule
\end{tabularx}
\end{center}

The identities of the transition algebra are unchanged across all four
examples.  Finite versus continuous spectrum, zero-mode imbalance, support
gap, and response measure are properties of the realization \(B\), not of
the abstract A/V relations alone.
% ===== END INLINED FILE: sections/09_examples.tex =====


% ===== BEGIN INLINED FILE: sections/10_scope_open_problems.tex =====
\section{Synthesis, realization data, and open extensions}
\label{sec:scope-open}

The preceding sections establish a local operator architecture.  This final
section states precisely what is common to all realizations, what depends on
the channel, and which extensions require additional mathematics.

\subsection{Architecture theorem}

\begin{theorem}[Universal transition-response architecture]
\label{thm:architecture}
Let \(B:\Dom B\subseteq\scrHp\to\scrHm\) be densely defined and closed.
Then the following structures are canonically determined.
\begin{enumerate}[label=(\alph*)]
  \item The nilpotent oriented pair \(\Aop_B=-d_B\),
  \(\Vop_B=d_B^*\), the positive partner Laplacian
  \(\DeltaB_B\), and the two odd branches \(\DB_B\), \(\KB_B\).
  \item The polar reflection \(\RB\), polar complex structure \(\IB\),
  support projection \(\SuppB\), and the supportwise rotation generated by
  \(\IB\).
  \item The two-term cohomology, paired nonzero partner spectra, and, when
  \(B\) is Fredholm, the index \(\ind B\).
\end{enumerate}
If a compact Hilbert-unitary symmetry action commuting with \(\fundJ\) is
supplied, Haar averaging gives a reducing invariant projection.  If
\(B=B_\mu\) arises from a norm-\(C^2\) projection field, its mixed returns
produce the Grassmann curvature.  If, in addition, the selected bundle is
smooth and finite rank, the same curvature produces the ordinary
Chern--Weil forms used here.  If the channel is
bounded and couples two self-adjoint diagonal generators, its
resolvent-dressed return produces the Schur--Feshbach effective operator.  If a self-adjoint realization is supplied
with a bounded port, its compressed resolvent produces a Herglotz response, a
Schur transfer function, and positive response kernels.
\end{theorem}

\begin{proof}
Items (a)--(c) collect
\cref{thm:AV-algebra},
\cref{thm:polar-properties,thm:polar-factorization},
\cref{thm:Hodge-decomposition,thm:partner-unitary-equivalence}, and
\cref{thm:index-formulas}.
The compact-symmetry statement is \cref{thm:Haar-projector}; the
variable-projection statement is
\cref{thm:tangent-realization,thm:projector-curvature,thm:chern-character};
the effective-dynamics statement is \cref{thm:Schur-factorization}; and the
response statement is
\cref{thm:response-Herglotz,thm:Schur-contractivity,thm:positive-Schur-kernel}.
\end{proof}

The theorem concerns the architecture generated by the channel and the
additional realization data.  It does not imply that different choices of
\(B\), \(\Pi\), \(H\), or \(\Gamma\) are physically or spectrally
equivalent.

\subsection{Realization datum}

A concrete model may be organized by the tuple
\begin{equation}
  \mathfrak R
  =\bigl(
    X,\scrK,\fundJ,\rho,\Pi_G,\Pi,B,\ket0,
    \mathbb L^{(\varepsilon)},H,\Gamma
   \bigr),
  \label{eq:realization-datum}
\end{equation}
with only the entries needed by that realization retained.  Their roles are:
\begin{center}
\small
\begin{tabularx}{\textwidth}{>{\raggedright\arraybackslash}p{0.22\textwidth} X}
\toprule
\textbf{Datum} & \textbf{Mathematical role} \\
\midrule
\(X\) & Parameter manifold for a moving projection, when present. \\
\((\scrK,\fundJ)\) & Bipartite Hilbert/Krein carrier and fixed grading. \\
\(\rho,\Pi_G\) & Optional compact symmetry action and its invariant Haar projection. \\
\(\Pi\) & Selected subbundle and Grassmann connection; finite rank is required for the ordinary Chern--Weil forms used here. \\
\(B\) & Closed transition channel; it determines the partner spectra. \\
\(\ket0\) & Optional normalized reference; it lies in \(\Dom B\) when the transition image is used, while ground-state or cyclicity properties require separate hypotheses. \\
\(\mathbb L^{(\varepsilon)}\) & Block generator for sector elimination. \\
\(H\) & Self-adjoint internal generator for passive response. \\
\(\Gamma\) & External port; it determines the compressed spectral measure. \\
\bottomrule
\end{tabularx}
\end{center}

A compact symmetry action may additionally supply the Haar projector
\(\Pi_G\) of \cref{thm:Haar-projector}.  It is a derived invariant
compression, not a replacement for the fixed grading or the moving projector.

A physical implementation must also specify a preparation procedure
for the selected port states and a measurement procedure yielding the
stated response.  Deriving an internal Hamiltonian or a source coupling
does not by itself derive those procedures.  A numerical or electronic
implementation can test its prescribed operations, precision and errors;
the physical identification of the underlying degrees of freedom is a
separate claim.

The data are not redundant.  In particular, the fixed fundamental symmetry
\(\fundJ\), a gauge-averaging projection, a moving physical projection, a
spectral half-line projection, and a boundary port are operators of different
types and should not share a symbol unless an explicit theorem identifies
them.

\subsection{Claim-status convention}

The formal environments already distinguish definitions, assumptions, and
theorems.  For proposed extensions beyond the present paper, the following
status language will be used.

\begin{center}
\small
\begin{tabularx}{\textwidth}{>{\raggedright\arraybackslash}p{0.18\textwidth} X}
\toprule
\textbf{Status} & \textbf{Meaning} \\
\midrule
Established & Proved in the present manuscript under the displayed hypotheses. \\
Model choice & A selected operator, domain, normalization, boundary condition, or realization datum. \\
Conditional & A theorem whose antecedent contains a separately named construction problem. \\
Open realization & A required map or object has not yet been constructed. \\
Obstruction & A known reason that a tempting identification does not follow from the current hypotheses. \\
\bottomrule
\end{tabularx}
\end{center}

The dependency rule is simple: an example cannot prove the abstract algebra;
a model choice cannot be cited as a derived normalization; and an open global
realization cannot be used to strengthen a local theorem.

\subsection{Established layers and boundaries}

\begin{center}
\small
\begin{tabularx}{\textwidth}{>{\raggedright\arraybackslash}p{0.28\textwidth} >{\raggedright\arraybackslash}p{0.18\textwidth} X}
\toprule
\textbf{Layer} & \textbf{Status} & \textbf{Boundary} \\
\midrule
Fixed \(\Aop/\Vop\) algebra & Established & Requires a densely defined closed channel and natural product domains. \\
Compact symmetry averaging & Established & Requires a compact group and a strongly continuous Hilbert-unitary representation; it yields an invariant projection, not a preferred dynamics. \\
Polar/cohomological structure & Established & Hodge decompositions without range closures require closed range; index formulas require Fredholmness. \\
Projector curvature & Established & Uses a norm-\(C^2\) orthogonal projection field; Chern forms here use a smooth finite-rank selected bundle. \\
Schur--Feshbach reduction & Established & The main statement assumes bounded channel coupling between self-adjoint diagonal blocks. \\
Passive response & Established & Full Herglotz positivity uses a Hilbert-self-adjoint realization and bounded port. \\
Heat-smoothed singular port & Established under the displayed Hilbert-scale hypothesis & Gaussian smoothing gives a bounded port for every \(\tau>0\); discrete series formulas additionally require the stated summability. \\
Positive-integer heat trace & Established example & Realized by \(B=a^*\); it does not determine a nontrivial-zero spectrum. \\
Symmetry-resolved response & Established under the displayed hypotheses & Ground preservation needs a no-crossing condition; the full defect is defined relative to the minimal carrier and the specified port. \\
Physical or arithmetic identification & Open realization & Requires additional operators, maps, domains, and comparison theorems. \\
\bottomrule
\end{tabularx}
\end{center}

\subsection{Open mathematical extensions}

The following problems arise naturally from the established core.

\begin{enumerate}[label=\textbf{G\arabic*.},leftmargin=3.5em]
  \item \textbf{Unbounded block coupling.}
  Extend the Schur--Feshbach factorization from bounded \(B\) to closed
  unbounded channels using form sums, boundary triples, or compatible graph
  norms while retaining precise resolvent and positivity statements.

  \item \textbf{\(C^2\) families of unbounded channels.}
  Develop a graph-topology version of the projector-tangent correspondence for
  parameter-dependent unbounded \(B_x\), including domain transport and
  differentiability of polar factors.

  \item \textbf{Families index.}
  Relate the curvature and Chern character of a family of active subspaces to
  the family of Fredholm channels and its index bundle.

  \item \textbf{Finite-to-continuum symmetry compression.}
  Control convergence of finite invariant projections, such as a cubic
  trivial-isotypic compression, to a continuum compact-group projection,
  together with convergence of the compressed generators and responses.

  \item \textbf{Singular boundary systems.}
  Replace the heat-smoothed bounded port by a full boundary-triple or rigged
  Hilbert-space treatment, and identify when the unsmoothed Weyl function is
  already well defined.

  \item \textbf{Krein-passive realization.}
  Determine conditions on the generator, port and pairing under which
  a Krein-self-adjoint block yields a positive physical response.
  Similarity to a Hilbert-self-adjoint operator does not by itself
  guarantee positivity of an unchanged Krein-paired port response.
  Any positive replacement pairing must be derived or stated as
  additional realization data.

  \item \textbf{Global assembly.}
  Specify a category of local channels and gluing morphisms for which closed
  excursions, indices, and response kernels assemble functorially over a
  global base.

  \item \textbf{Arithmetic realization.}
  The oscillator channel gives the exact positive-integer heat trace and its
  Mellin transform.  A further construction would have to produce arithmetic
  primitive data, an appropriate global determinant, and a positive
  comparison kernel without inserting those objects by definition.
\end{enumerate}

\subsection{Visible limitations}

The canonical realizations make several boundaries explicit.
\begin{enumerate}[label=\textbf{O\arabic*.},leftmargin=3.5em]
  \item The translation channel has continuous spectrum touching zero and no
  support gap.  The transition algebra alone therefore does not force a gap.
  \item Krein self-adjointness does not by itself imply passivity of the full
  compressed resolvent; the response theorem uses a Hilbert-self-adjoint
  realization.
  \item Compact-group averaging produces an invariant subspace but does not
  determine a Hamiltonian, action, or preferred structure group.
  \item Projector curvature produces characteristic classes, whereas the
  holonomy of an arbitrary closed loop need not be quantized.
  \item The oscillator integer heat trace does not construct a spectrum of
  nontrivial zeta zeros or the endpoint factors of the entire completed
  function.
\end{enumerate}

\subsection{Conclusion}

The common content is the functorial passage from a closed channel to its
oriented differential pair, positive return, polar support geometry, and
cohomology; from a moving channel to curvature; from a resolvent-dressed
channel to effective dynamics; and from a ported self-adjoint realization to
a passive positive response.  Spectral type, topology, and physical content
remain properties of the realization datum rather than consequences of the
abstract orientation algebra alone.
% ===== END INLINED FILE: sections/10_scope_open_problems.tex =====


\clearpage
\appendix

% ===== BEGIN INLINED FILE: appendices/A_domains.tex =====
\section{Domain bookkeeping}
\label{app:domains}

This appendix collects the natural domains of every unbounded operator used
in the main text. No block product is intended outside its composition
domain.

\subsection{First-order operators}

For the closed densely defined channel
\(B:\Dom B\subseteq\scrHp\to\scrHm\),
\begin{align}
  \Dom d_B
  &=\Dom\Aop_B
  =\Dom B\oplus\scrHm,\\
  \Dom d_B^*
  &=\Dom\Vop_B
  =\scrHp\oplus\Dom B^*,\\
  \Dom\DB_B
  &=\Dom\KB_B
  =\Dom B\oplus\Dom B^*.
\end{align}
The square domains of the nilpotent one-way maps are unchanged:
\begin{equation}
  \Dom(\Aop_B^2)=\Dom\Aop_B,
  \qquad
  \Dom(\Vop_B^2)=\Dom\Vop_B.
\end{equation}

\subsection{Ordered products and Laplacian}

The ordered products have domains
\begin{align}
  \Dom(\Vop_B\Aop_B)
  &=\Dom(B^*B)\oplus\scrHm,\\
  \Dom(\Aop_B\Vop_B)
  &=\scrHp\oplus\Dom(BB^*).
\end{align}
Indeed, for \((x,y)\in\Dom\Aop_B\),
\(\Aop_B(x,y)=(0,-Bx)\) lies in \(\Dom\Vop_B\) precisely when
\(Bx\in\Dom B^*\), equivalently \(x\in\Dom(B^*B)\). The other product is
symmetric.

The common domain on which the anticommutator, commutator, and squares of the
odd branches are evaluated is
\begin{equation}
  \Dom\DeltaB_B
  =\Dom(B^*B)\oplus\Dom(BB^*).
\end{equation}
On this domain,
\begin{equation}
  \DB_B^2=\DeltaB_B,
  \qquad
  \KB_B^2=-\DeltaB_B,
  \qquad
  \{\Aop_B,\Vop_B\}=-\DeltaB_B.
\end{equation}

\subsection{Reference-state relations}

The definition \(\ket\Omega=\Aop_B\ket0\) requires only
\(\ket0\in\Dom B\). The nilpotent identity
\(\Aop_B\ket\Omega=0\) then holds automatically because every vector in the
negative sector belongs to \(\Dom\Aop_B\). The return expression
\(\Vop_B\ket\Omega=-B^*B\ket0\) requires the stronger condition
\(\ket0\in\Dom(B^*B)\).

\subsection{Polar factors}

The partial isometry \(u_B\), support projections
\(\pi_{+,B},\pi_{-,B},\Pi_B\), polar reflection \(\mathcal R_B\), polar
complex structure \(\mathcal I_B\), and normalized transitions
\(\widehat{\Aop}_B,\widehat{\Vop}_B\) are bounded and defined on all of
\(\scrK\). The square-root domain is
\begin{equation}
  \Dom\DeltaB_B^{1/2}
  =\Dom B\oplus\Dom B^*
  =\Dom\DB_B
  =\Dom\KB_B.
\end{equation}
Consequently, the factorizations
\begin{equation}
  \DB_B=\mathcal R_B\DeltaB_B^{1/2},
  \qquad
  \KB_B=\mathcal I_B\DeltaB_B^{1/2}
\end{equation}
have exactly this first-order domain. The reverse-order products have the
same domain because the polar factors intertwine \(|B|\) and \(|B^*|\).

Let \(\DeltaB_{B,0}^{-1/2}\) denote the generalized inverse defined by
spectral calculus from the function \(\lambda\mapsto\lambda^{-1/2}\) for
\(\lambda>0\), with value zero at \(\lambda=0\). Its domain is
\begin{equation}
  \Dom\DeltaB_{B,0}^{-1/2}
  =\left\{\xi\in\scrK:
  \int_{(0,\infty)}\lambda^{-1}\,d\mu_\xi^{\DeltaB_B}(\lambda)<\infty
  \right\},
\end{equation}
where \(\mu_\xi^{\DeltaB_B}\) is the scalar spectral measure of
\(\DeltaB_B\) at \(\xi\). On \(\Ran\DeltaB_B^{1/2}\), these products agree with the restrictions of
\(\mathcal R_B\) and \(\mathcal I_B\) to the active support.  They extend
uniquely there; extension by zero on \(\Ker\DeltaB_B\) gives the bounded
polar operators used in the main text.

\subsection{Partner restrictions and cohomology}

The active restrictions of \(B^*B\) and \(BB^*\) are self-adjoint on the
reducing subspaces
\begin{equation}
  (\Ker B)^\perp,
  \qquad
  (\Ker B^*)^\perp.
\end{equation}
The two-term cohomology uses \(B\) with its original closed domain
\(\Dom B\); no bounded extension of \(B\) is implied. If \(\Ran B\) is not
closed, the ordinary quotient \(\scrHm/\Ran B\) may be non-Hausdorff, while
\(\scrHm/\overline{\Ran B}\) is the reduced Hilbert-space cohomology.

The heat-supertrace identity is asserted only when the two heat semigroups
are trace class. This is an additional spectral hypothesis and does not
follow from Fredholmness alone.

\subsection{Moving projections}

The projection field \(\Pi:X\to\mathcal B(\scrK)\) is assumed to be
\(C^2\) in the operator norm.  Consequently, every tangent
\(\partial_\mu\Pi\) and every channel
\(B_\mu=\Pi^\perp(\partial_\mu\Pi)\Pi\) is bounded and defined on the full
fiber.  All operators
\begin{equation}
  \Aop_\mu,\quad
  \Vop_\mu,\quad
  \DB_{B_\mu},\quad
  \KB_{B_\mu}
\end{equation}
are therefore bounded on the moving decomposition
\(\mathscr E_x\oplus\mathscr E_x^\perp\).  The curvature products
\(B_\mu^*B_\nu\) require no additional domain hypotheses.

\subsection{Schur--Feshbach block generators}

In \cref{sec:effective-dynamics}, the coupling channel is assumed bounded:
\(B\in\mathcal B(\scrHp,\scrHm)\).  The block generators
\begin{equation}
  \mathbb L^{(\varepsilon)}
  =\begin{pmatrix}L_+&B^*\\ \varepsilon B&L_-\end{pmatrix}
\end{equation}
have the common domain
\begin{equation}
  \Dom\mathbb L^{(\varepsilon)}
  =\Dom L_+\oplus\Dom L_-.
\end{equation}
For \(z\in\rho(L_-)\), the return
\(\Sigma_B(z)=B^*(L_--z)^{-1}B\) is bounded on \(\scrHp\), and the Schur
pencil
\begin{equation}
  S_\varepsilon(z)=L_+-z-\varepsilon\Sigma_B(z)
\end{equation}
is closed on \(\Dom L_+\).  The block inverse formula is asserted only when
\(0\in\rho(S_\varepsilon(z))\).  No claim about an unbounded off-diagonal
coupling is made in this pass.

\subsection{Boundary response}

In \cref{sec:boundary-response}, the response generator
\(H:\Dom H\subseteq\scrX\to\scrX\) is self-adjoint and the port
\(\Gamma\in\mathcal B(\scrE,\scrX)\) is bounded.  Hence
\((H-z)^{-1}\Gamma\) is bounded for every \(z\in\CC\setminus\RR\), and
\begin{equation}
  M_\Gamma(z)=\Gamma^*(H-z)^{-1}\Gamma
\end{equation}
is defined on all of \(\scrE\).  The Cayley denominator
\(M_\Gamma(z)+i\mathbf1_{\scrE}\) is boundedly invertible in the upper
half-plane by \cref{lem:Cayley-denominator}.

For \(s\ge0\), the Hilbert scale is
\begin{equation}
  \scrX_s=\Dom(\mathbf1+H^2)^{s/2},
\end{equation}
with anti-dual \(\scrX_{-s}\).  If
\(\gamma:\scrX_s\to\scrE\) is bounded, then
\(\gamma^*:\scrE\to\scrX_{-s}\), and functional calculus gives a bounded map
\begin{equation}
  e^{-\tau H^2/2}:\scrX_{-s}\longrightarrow\scrX
\end{equation}
for every \(\tau>0\).  Thus
\(\Gamma_{\tau,\gamma}=e^{-\tau H^2/2}\gamma^*\) is an ordinary bounded
port.  The internal operator \(H\) and its domain are unchanged by this
smoothing.

\subsection{Example domains}

For the translation channel,
\begin{equation}
  D_x=-i\frac{d}{dx},
  \qquad
  \Dom D_x=H^1(\RR),
  \qquad
  \Dom D_x^2=H^2(\RR).
\end{equation}
For the oscillator channel, \(a\) and \(a^*\) have their standard maximal
closed domains in \(\ell^2(\NN_0)\); the products \(aa^*=N+\mathbf1\) and
\(a^*a=N\) are evaluated on the corresponding number-operator domains.
For the quasiperiodic compact channel,
\begin{equation}
  D_{L,\vartheta}=-i\frac{d}{du},
\end{equation}
with domain
\begin{equation}
  \Dom D_{L,\vartheta}
  =\{f\in H^1(-L,L):f(L)=e^{i\vartheta}f(-L)\}.
\end{equation}
Point evaluation at \(u_0\) is used only as a bounded functional on
\(H^s(-L,L)\) with \(s>1/2\); its delta adjoint is smoothed before entering
the \(L^2\) response.

\subsection{Reference-state geometry}

For \cref{thm:reference-fidelity}, the vector \(\psi\) lies in
\(\Dom H_+\), so the expectation of \(H_+\) and the orthogonal decomposition
are well defined.  Its gap corollary is stated on the form domain
\(\Dom((H_+-E_0)^{1/2})\), where
\(E_{\mathrm{exc}}(\psi)=\|(H_+-E_0)^{1/2}\psi\|^2\).  The rank-one projection \(P_0\) is bounded.
No transition-channel domain condition is needed unless \(H_+\) is later
identified with \(B^*B\) or another channel-derived operator.

\subsection{Characteristic classes}

The Chern--Weil section assumes that \(\Pi(x)\) has finite rank and is smooth
enough for the operator-valued differential forms displayed there.  The
fiber trace is then finite.  For infinite-rank bundles, trace-class or
regularized characteristic forms require additional hypotheses and are not
covered by \cref{thm:chern-character}.
% ===== END INLINED FILE: appendices/A_domains.tex =====

\clearpage

% ===== BEGIN INLINED FILE: appendices/B_notation.tex =====
\section{Notation index}
\label{app:notation}

\subsection{Fixed-channel algebra}

\begin{center}
\small
\begin{tabularx}{\textwidth}{>{\raggedright\arraybackslash}p{0.24\textwidth} >{\raggedright\arraybackslash}p{0.28\textwidth} X}
\toprule
\textbf{Symbol} & \textbf{Type} & \textbf{Definition} \\
\midrule
\(\scrHp,\scrHm\) & Hilbert spaces & Positive and negative carrier spaces. \\
\(\scrK\) & Hilbert/Krein space & \(\scrHp\oplus\scrHm\) with fundamental decomposition \(\scrHp\dsum\scrHm\). \\
\(\fundJ\) & Bounded involution & \(\diag(\idop_+,-\idop_-)\). \\
\(\Pp,\Pm\) & Orthogonal projections & Fixed projections onto \(\scrHp\) and \(\scrHm\). \\
\(X^*\) & Hilbert adjoint & Ordinary adjoint in the positive Hilbert topology. \\
\(X^\sharp\) & Krein adjoint & \(\fundJ X^*\fundJ\). \\
\(B\) & Closed channel & \(B:\Dom B\subseteq\scrHp\to\scrHm\). \\
\(d_B\) & Differential & \(\bigl(\begin{smallmatrix}0&0\\B&0\end{smallmatrix}\bigr)\). \\
\(d_B^*\) & Codifferential & \(\bigl(\begin{smallmatrix}0&B^*\\0&0\end{smallmatrix}\bigr)\). \\
\(\Aop_B\) & Oriented transition & \(-d_B\). \\
\(\Vop_B\) & Oriented transition & \(d_B^*\). \\
\(\DB_B\) & Hilbert/Hodge branch & \(\Vop_B-\Aop_B=d_B+d_B^*\). \\
\(\KB_B\) & Krein/rotational branch & \(\Vop_B+\Aop_B=d_B^*-d_B\). \\
\(\DeltaB_B\) & Positive Laplacian & \(\diag(B^*B,BB^*)\). \\
\(\mathscr C_{+,B},\mathscr C_{-,B}\) & Ordered cycles & \(-\Vop_B\Aop_B\) and \(-\Aop_B\Vop_B\). \\
\(u_B\) & Polar partial isometry & Polar phase in \(B=u_B|B|\). \\
\(\pi_{+,B},\pi_{-,B}\) & Support projections & \(u_B^*u_B\) and \(u_Bu_B^*\). \\
\(\Pi_B\) & Active-support projection & \(\diag(\pi_{+,B},\pi_{-,B})=\mathbf1_{(0,\infty)}(\DeltaB_B)\). \\
\(\mathcal R_B\) & Polar reflection & \(\bigl(\begin{smallmatrix}0&u_B^*\\u_B&0\end{smallmatrix}\bigr)\), with \(\mathcal R_B^2=\Pi_B\). \\
\(\mathcal I_B\) & Polar complex structure & \(\bigl(\begin{smallmatrix}0&u_B^*\\-u_B&0\end{smallmatrix}\bigr)\), with \(\mathcal I_B^2=-\Pi_B\). \\
\(\widehat{\Aop}_B,\widehat{\Vop}_B\) & Normalized transitions & The oriented pair with \(B\) replaced by \(u_B\). \\
\(H_B^0,H_B^1\) & Cohomology & \(\Ker B\) and \(\scrHm/\Ran B\). \\
\(\overline H_B^1\) & Reduced cohomology & \(\scrHm/\overline{\Ran B}\). \\
\(\ind B\) & Fredholm index & \(\dim\Ker B-\dim\Ker B^*\). \\
\(\operatorname{Str}_{\fundJ}\) & Graded trace & \(\Tr(\fundJ\,\cdot)\) when defined. \\
\(\ket0\) & Reference vector & A normalized vector in \(\Dom B\subseteq\scrHp\). \\
\(\ket\Omega\) & Transition image & \(\Aop_B\ket0\in\scrHm\). \\
\bottomrule
\end{tabularx}
\end{center}

\subsection{Moving projections and effective dynamics}

\begin{center}
\small
\begin{tabularx}{\textwidth}{>{\raggedright\arraybackslash}p{0.24\textwidth} >{\raggedright\arraybackslash}p{0.28\textwidth} X}
\toprule
\textbf{Symbol} & \textbf{Type} & \textbf{Definition} \\
\midrule
\(\Pi_G\) & Symmetry-average projection & Haar projection onto the invariant subspace of a compact unitary action. \\
\(\Pi(x),\Pi^\perp(x)\) & Moving projections & Smooth orthogonal projection and its complement. \\
\(\mathfrak J_\Pi\) & Moving fundamental symmetry & \(2\Pi-\mathbf1\). \\
\(B_\mu\) & Tangent channel & \(\Pi^\perp(\partial_\mu\Pi)\Pi\). \\
\(\mathsf A_\mu,\mathsf V_\mu\) & Tangent transitions & Oriented pair associated with \(B_\mu\). \\
\(\mathcal A_\mu^{(W)}\) & Local connection & \(W^*\partial_\mu W\), in the anti-Hermitian convention. \\
\(\mathcal F_{\mu\nu}^{\Pi}\) & Projector curvature & \(\Pi[\partial_\mu\Pi,\partial_\nu\Pi]\Pi\). \\
\(\mathscr C_{\mu\nu}^{\Pi}\) & Mixed return & \(B_\mu^*B_\nu=-\mathsf V_\mu\mathsf A_\nu\) on \(\operatorname{Ran}\Pi\). \\
\(\mathbb L^{(\varepsilon)}\) & Block generator & \(\bigl(\begin{smallmatrix}L_+&B^*\\\varepsilon B&L_-\end{smallmatrix}\bigr)\). \\
\(R_-(z)\) & Complementary resolvent & \((L_--z)^{-1}\). \\
\(\Sigma_B(z)\) & Self-energy/return & \(B^*R_-(z)B\). \\
\(L_{\mathrm{eff}}^{(\varepsilon)}(z)\) & Effective operator & \(L_+-\varepsilon\Sigma_B(z)\). \\
\(S_\varepsilon(z)\) & Schur pencil & \(L_{\mathrm{eff}}^{(\varepsilon)}(z)-z\). \\
\(\rho(g)\) & Group representation & Reserved for a symmetry-group action. \\
\(U(\phi)\) & Scalar potential & Reserved for field-theoretic realizations. \\
\bottomrule
\end{tabularx}
\end{center}

\subsection{Boundary response and realizations}

\begin{center}
\small
\begin{tabularx}{\textwidth}{>{\raggedright\arraybackslash}p{0.24\textwidth} >{\raggedright\arraybackslash}p{0.28\textwidth} X}
\toprule
\textbf{Symbol} & \textbf{Type} & \textbf{Definition} \\
\midrule
\(\scrX\) & Internal Hilbert space & Carrier of the self-adjoint response generator. \\
\(\scrE\) & Port Hilbert space & External input--output space. \\
\(H\) & Response generator & Self-adjoint operator on \(\scrX\). \\
\(\Gamma\) & Input map & Bounded map \(\scrE\to\scrX\). \\
\(\Sigma_\Gamma\) & Compressed spectral measure & \(\Gamma^*E_H(\cdot)\Gamma\). \\
\(M_\Gamma(z)\) & Herglotz response & \(\Gamma^*(H-z)^{-1}\Gamma\). \\
\(C_\Gamma(t)\) & Time response & \(\Gamma^*e^{-itH}\Gamma\). \\
\(\mathcal K_M(z,w)\) & Herglotz kernel & \((M_\Gamma(z)-M_\Gamma(w)^*)/(z-\bar w)\). \\
\(\mathscr S_\Gamma(z)\) & Schur response & Cayley transform of \(M_\Gamma(z)\). \\
\(\mathcal K_{\mathscr S}(z,w)\) & de Branges--Rovnyak kernel & \((\mathbf1-\mathscr S_\Gamma(z)\mathscr S_\Gamma(w)^*)/[-i(z-\bar w)]\). \\
\(\mathscr X_{\min}(H,\Gamma)\) & Reachable subspace & Closed span of \((H-z)^{-1}\Gamma e\). \\
\(\scrX_{\pm s}(H)\) & Hilbert scale & Scale generated by \((\mathbf1+H^2)^{s/2}\). \\
\(\Gamma_{\tau,\gamma}\) & Smoothed port & \(e^{-\tau H^2/2}\gamma^*\). \\
\(D_{L,\vartheta}\) & Compact first-order generator & \(-i\,d/du\) with quasiperiodic boundary condition. \\
\bottomrule
\end{tabularx}
\end{center}

\subsection{Reference-state and topological notation}

\begin{center}
\small
\begin{tabularx}{\textwidth}{>{\raggedright\arraybackslash}p{0.24\textwidth} >{\raggedright\arraybackslash}p{0.28\textwidth} X}
\toprule
\textbf{Symbol} & \textbf{Type} & \textbf{Definition} \\
\midrule
\(P_0\) & Rank-one projection & \(\ket0\bra0\) when \(\ket0\) is a normalized ground state. \\
\(Q_0\) & Complementary projection & \(\mathbf1-P_0\). \\
\(v_\psi\) & Reference fidelity & \(|\braket0\psi|^2\). \\
\(\eta_\psi\) & Conditional excitation vector & \(Q_0\psi/\sqrt{1-v_\psi}\). \\
\(Z_\psi\) & Conditional excitation scale & \(\inner{\eta_\psi}{(H_+-E_0)\eta_\psi}\). \\
\(\operatorname{ch}_k(\mathscr E)\) & Characteristic form & Degree-\(2k\) component generated from \(\mathcal F^\Pi\). \\
\(c_1(\mathscr E)\) & Integral characteristic class & First Chern class of a finite-rank selected bundle. \\
\(\mathfrak R\) & Realization datum & Tuple collecting only the carrier, channel, projection, generator, and port data used by a concrete model. \\
\bottomrule
\end{tabularx}
\end{center}
% ===== END INLINED FILE: appendices/B_notation.tex =====


\clearpage
\begin{thebibliography}{99}
\bibitem{AzizovIokhvidov}
T.~Ya. Azizov and I.~S. Iokhvidov,
\emph{Linear Operators in Spaces with an Indefinite Metric},
Wiley, 1989.

\bibitem{Bognar}
J.~Bogn\'ar,
\emph{Indefinite Inner Product Spaces},
Springer, 1974.

\bibitem{Kato}
T.~Kato,
\emph{Perturbation Theory for Linear Operators},
Springer, 2nd ed., 1995.

\bibitem{ReedSimon}
M.~Reed and B.~Simon,
\emph{Methods of Modern Mathematical Physics, Vol. I: Functional Analysis},
Academic Press, 1980.

\bibitem{BruningLesch}
J.~Br\"uning and M.~Lesch,
``Hilbert complexes,''
\emph{Journal of Functional Analysis} \textbf{108} (1992), 88--132,
\href{https://doi.org/10.1016/0022-1236(92)90147-B}{doi:10.1016/0022-1236(92)90147-B}.

\bibitem{Donoghue}
W.~F. Donoghue, Jr.,
\emph{Monotone Matrix Functions and Analytic Continuation},
Springer, 1974.

\bibitem{Aronszajn}
N.~Aronszajn,
``Theory of reproducing kernels,''
\emph{Transactions of the American Mathematical Society} \textbf{68}
(1950), 337--404,
\href{https://doi.org/10.1090/S0002-9947-1950-0051437-7}{doi:10.1090/S0002-9947-1950-0051437-7}.

\bibitem{deBrangesRovnyak}
L.~de Branges and J.~Rovnyak,
\emph{Square Summable Power Series},
Holt, Rinehart and Winston, 1966.


\bibitem{KobayashiNomizu}
S.~Kobayashi and K.~Nomizu,
\emph{Foundations of Differential Geometry, Vol. II},
Wiley, 1969.

\bibitem{Husemoller}
D.~Husemoller,
\emph{Fibre Bundles},
Springer, 3rd ed., 1994.
\bibitem{Gilbert1963}
E.~G. Gilbert,
``Controllability and observability in multivariable control systems,''
\emph{Journal of the Society for Industrial and Applied Mathematics,
Series A: Control} \textbf{1}(2) (1963), 128--151,
\href{https://doi.org/10.1137/0301009}{doi:10.1137/0301009}.

\bibitem{Kalman1963}
R.~E. Kalman,
``Mathematical description of linear dynamical systems,''
\emph{Journal of the Society for Industrial and Applied Mathematics,
Series A: Control} \textbf{1}(2) (1963), 152--192,
\href{https://doi.org/10.1137/0301010}{doi:10.1137/0301010}.

\end{thebibliography}

\end{document}
